\documentclass[11pt,a4paper]{article}

\usepackage[T1]{fontenc}
\usepackage[utf8]{inputenc}
\usepackage{newtxtext}
\usepackage[a4paper,margin=27mm,headheight=15pt]{geometry}
\usepackage{microtype}
\usepackage{mathtools,amssymb,amsthm}
\usepackage{newtxmath}
\usepackage{bm}
\usepackage{booktabs,longtable,array,multirow}
\usepackage{enumitem}
\usepackage{xcolor}
\usepackage[most]{tcolorbox}
\usepackage{fancyhdr}
\usepackage{xurl}
\usepackage{hyperref}
\usepackage[nameinlink,noabbrev]{cleveref}
\usepackage{listings}
\usepackage{needspace}
\usepackage{csquotes}

\hypersetup{
  colorlinks=true,
  linkcolor=blue!45!black,
  citecolor=green!35!black,
  urlcolor=blue!55!black,
  pdftitle={Arithmetic Rademacher Transseries for the Partition Numbers and Their Asymptotic Inverse},
  pdfauthor={Prepared for the ProveIt project},
  pdfsubject={Partition asymptotics, transseries, Lambert W inversion, coefficient calculus},
  pdfkeywords={partition function, Rademacher series, asymptotic transseries, inverse function, Lambert W, Bell polynomials, Stirling numbers}
}

\pagestyle{fancy}
\fancyhf{}
\fancyhead[L]{\small Partition-number transseries}
\fancyhead[R]{\small \nouppercase{\leftmark}}
\fancyfoot[C]{\thepage}

\setlength{\parindent}{1.2em}
\setlength{\parskip}{0.25em}
\setlist{nosep,leftmargin=2em}
\allowdisplaybreaks
\numberwithin{equation}{section}

\newtheorem{theorem}{Theorem}[section]
\newtheorem{proposition}[theorem]{Proposition}
\newtheorem{lemma}[theorem]{Lemma}
\newtheorem{corollary}[theorem]{Corollary}
\newtheorem{definition}[theorem]{Definition}
\newtheorem{algorithm}[theorem]{Algorithm}
\newtheorem{conjecture}[theorem]{Conjecture}
\theoremstyle{remark}
\newtheorem{remark}[theorem]{Remark}
\newtheorem{example}[theorem]{Example}

\newtcolorbox{resultbox}[1][]{
  enhanced,breakable,colback=blue!2,colframe=blue!45!black,
  boxrule=0.5pt,arc=1mm,left=1.5mm,right=1.5mm,top=1mm,bottom=1mm,
  title={#1},fonttitle=\bfseries
}
\newtcolorbox{warningbox}[1][]{
  enhanced,breakable,colback=orange!3,colframe=orange!55!black,
  boxrule=0.5pt,arc=1mm,left=1.5mm,right=1.5mm,top=1mm,bottom=1mm,
  title={#1},fonttitle=\bfseries
}

\newcommand{\e}{\mathrm e}
\newcommand{\ii}{\mathrm i}
\newcommand{\dd}{\,\mathrm d}
\newcommand{\CC}{\mathbb C}
\newcommand{\RR}{\mathbb R}
\newcommand{\QQ}{\mathbb Q}
\newcommand{\NN}{\mathbb N}
\newcommand{\ZZ}{\mathbb Z}
\newcommand{\What}{\widehat{\ZZ}}
\newcommand{\coef}[2]{[#1^{#2}]}
\newcommand{\Ccal}{\mathcal C}
\newcommand{\Ecal}{\mathcal E}
\newcommand{\Pcal}{\mathcal P}
\newcommand{\Rcal}{\mathcal R}
\newcommand{\Bcal}{\mathcal B}
\newcommand{\smallO}{\mathrm o}
\newcommand{\bigO}{\mathrm O}
\newcommand{\rising}[2]{(#1)_{#2}}

\title{\bfseries Arithmetic Rademacher Transseries for the Partition Numbers\\[0.35em]
\large Exact Exponential Sectors, All-Order Coefficients, and Asymptotic Inversion}
\author{Prepared for the ProveIt research program}
\date{3 September 2026}

\begin{document}
\maketitle

\begin{abstract}
Let $p(n)$ be the ordinary partition number.  The familiar Hardy--Ramanujan term captures only the first level of a much richer asymptotic object.  Starting from Rademacher's convergent series, we introduce the shifted phase
\[
 N=n-\frac1{24},\qquad x=\pi\sqrt{\frac{2N}{3}},\qquad
 \Ccal=\frac{\pi^2}{6\sqrt3},
\]
and prove the exact sector decomposition
\[
 p(n)=\frac{\Ccal}{x^2}\sum_{k\ge1}\frac{A_k(n)}{\sqrt{k}}
 \left\{\e^{x/k}\left(1-\frac{k}{x}\right)
       +\e^{-x/k}\left(1+\frac{k}{x}\right)\right\}.
\]
Thus the full forward transseries is not guessed from a divergent power series: it is the reorganized, convergent Rademacher expansion itself.  Each sector has only a linear amplitude in $x^{-1}$ in the natural shifted variable.  Re-expansion in $n^{-1/2}$ produces an infinite coefficient sequence.  We give (i) a coefficient-generating function valid for every $k$ and both signs, (ii) a complete Bell-polynomial formula and a triangular recurrence for all coefficients, and (iii) a proof of the finite O'Sullivan formula for the dominant coefficients by a Catalan change of variables and Lagrange inversion.

The inverse problem has two layers.  The elementary carrier $y=\Ccal\e^X/X^2$ is inverted exactly by
\[
 X=-2W_{-1}\!\left(-\frac12\sqrt{\frac{\Ccal}{y}}\right).
\]
Its complete iterated-logarithm expansion has coefficient polynomials
\[
 P_j(L)=2^{j+1}\sum_{q=1}^{j}\frac{s(j,q)}{(j+1-q)!}L^{j+1-q},
 \qquad L=\log\log(y/\Ccal),
\]
where $s(j,q)$ are signed Stirling numbers of the first kind.  Inverting the full dominant Rademacher sector adds a convergent rational series $U(1/X)=\sum b_mX^{-m}$; we derive a coefficient recurrence and list its first terms.  Finally, because $A_k(n)$ is periodic modulo $k$, the exponentially small inverse sectors form an arithmetic family.  At every finite Rademacher level $K$, fixing $n$ modulo $\operatorname{lcm}(1,\ldots,K)$ turns the problem into an ordinary analytic transseries inversion.  We give both a multivariate Newton--Hensel recurrence for every sector coefficient and a closed Lagrange--Buer\-mann/Bell-polynomial formula.  Compatible residue classes yield a natural profinite phase parameter in $\widehat{\mathbb Z}$.

We distinguish this smooth/asymptotic inverse from the integer-valued generalized inverse of the discrete sequence.  A final rounding layer is unavoidable for arbitrary real ordinates, while at exact ordinates $y=p(n)$ the dominant-sector inverse already differs from $n$ by an exponentially small quantity and therefore recovers $n$ by nearest-integer rounding for all sufficiently large $n$.  For completeness, we also give the full transseries of the multiplicative reciprocal $1/p(n)$.
\end{abstract}

\tableofcontents

\section{Introduction and statement of the program}

The partition function $p(n)$ is defined by
\begin{equation}\label{eq:partition-gf}
 \sum_{n\ge0}p(n)q^n=\prod_{m\ge1}\frac1{1-q^m},\qquad |q|<1.
\end{equation}
It is one of the central coefficient sequences in additive number theory.  Hardy and Ramanujan proved the celebrated leading approximation
\begin{equation}\label{eq:HR-leading}
 p(n)\sim \frac{1}{4\sqrt3\,n}
 \exp\!\left(\pi\sqrt{\frac{2n}{3}}\right),
\end{equation}
and Rademacher replaced the asymptotic formula by an exact convergent series involving Dedekind sums.  Modern work gives all algebraic orders and explicit error bounds; see Hardy--Ramanujan~\cite{HardyRamanujan1918}, Rademacher~\cite{Rademacher1937,Rademacher1938,Rademacher1943}, Lehmer~\cite{Lehmer1937,Lehmer1938,Lehmer1939}, O'Sullivan~\cite{OSullivan2023}, and Banerjee--Paule--Radu--Schneider~\cite{BanerjeeEtAl2022}.  The exact Rademacher formula and its leading asymptotic were also formalized in Isabelle/HOL in 2026~\cite{EberlPaulson2026}.

The purpose of this article is not to reprove the circle method.  We take Rademacher's theorem as the analytic input and develop, from it, a complete coefficient calculus for the forward and inverse transseries.  This choice has three advantages.

\begin{enumerate}
\item The complete exponential hierarchy is present from the beginning.  There is no need to infer exponentially small sectors from a divergent perturbative series.
\item The natural shifted phase $x=\pi\sqrt{2(n-1/24)/3}$ turns every Rademacher summand into two exponentials with elementary amplitudes.
\item All inverse coefficients can be generated by triangular recurrences, Bell polynomials, and Lagrange--Buer\-mann formulas of the kind developed in the accompanying coefficient-calculus and transseries documents of the ProveIt repository~\cite{ProveItCCC,ProveItTransseries}.
\end{enumerate}

There is an important semantic issue.  The phrase ``inverse of the partition numbers'' can mean either the reciprocal $1/p(n)$ or the inverse function that recovers $n$ from a large value of $p(n)$.  The main subject of this article is the latter.  Because $p$ is a sequence, however, an inverse function is not canonical until one chooses an interpolation or restricts to the exact range $\{p(n)\}$.  We make this distinction precise in \cref{sec:meaning-inverse}.  The multiplicative reciprocal is simpler and is treated in \cref{sec:reciprocal}.

\subsection{The main formulas at a glance}

Set
\begin{equation}\label{eq:master-variables}
 a=\frac1{24},\qquad N=n-a,\qquad
 \lambda=\pi\sqrt{\frac23},\qquad x=\lambda\sqrt N,
 \qquad \Ccal=\frac{\pi^2}{6\sqrt3}.
\end{equation}
Then the exact forward transseries is
\begin{equation}\label{eq:intro-exact}
 p(n)=\frac{\Ccal}{x^2}\sum_{k\ge1}\frac{A_k(n)}{\sqrt{k}}
 \left[\e^{x/k}\left(1-\frac{k}{x}\right)
       +\e^{-x/k}\left(1+\frac{k}{x}\right)\right].
\end{equation}
In unshifted powers $t=n^{-1/2}$, the $(k,\sigma)$ sector, $\sigma\in\{+1,-1\}$, is
\begin{equation}\label{eq:intro-sector}
 \frac{A_k(n)}{4\sqrt{3k}\,n}\,
 \e^{\sigma\lambda\sqrt n/k}
 \sum_{m\ge0}c_{k,\sigma,m}n^{-m/2},
\end{equation}
where
\begin{align}
 \sum_{m\ge0}c_{k,\sigma,m}t^m
 &=(1-at^2)^{-1}
 \exp\!\left[\frac{\sigma\lambda}{kt}
       \left(\sqrt{1-at^2}-1\right)\right]
 \left(1-\frac{\sigma kt}{\lambda\sqrt{1-at^2}}\right).
 \label{eq:intro-coeff-gf}
\end{align}
This one formula generates every algebraic coefficient in every Rademacher sector.

For the inverse, let
\begin{equation}\label{eq:intro-X}
 R=\log\frac{y}{\Ccal},\qquad
 X=-2W_{-1}\!\left(-\frac12\sqrt{\frac{\Ccal}{y}}\right).
\end{equation}
Then $X-2\log X=R$.  If $L=\log R$, the complete polynomial--logarithmic expansion is
\begin{equation}\label{eq:intro-Xexp}
 X\sim R+2L+\sum_{j\ge1}\frac{P_j(L)}{R^j},\qquad
 P_j(L)=2^{j+1}\sum_{q=1}^{j}
 \frac{s(j,q)}{(j+1-q)!}L^{j+1-q}.
\end{equation}
The inverse of the full dominant positive sector is
\begin{equation}\label{eq:intro-x0}
 x_0=X+U(X^{-1}),\qquad U(z)=\sum_{m\ge1}b_mz^m,
\end{equation}
where $U$ is the unique formal solution of
\begin{equation}\label{eq:intro-Ueq}
 U+\log(1-z+zU)-3\log(1+zU)=0.
\end{equation}
The complete arithmetic exponential correction is then obtained by either of two equivalent general formulas:
\begin{itemize}
\item the coefficient-by-coefficient recurrence in \cref{thm:multivariate-recurrence};
\item the Lagrange--Buer\-mann composition formula in \cref{thm:LB-full-inverse}.
\end{itemize}
Finally,
\begin{equation}\label{eq:intro-n-from-x}
 n=a+\frac{x^2}{\lambda^2}=\frac1{24}+\frac{3x^2}{2\pi^2}.
\end{equation}

\section{Notation, Dedekind sums, and coefficient polynomials}\label{sec:notation}

\subsection{Dedekind sums and Rademacher coefficients}

For a real number $u$, define the sawtooth function
\begin{equation}
 ((u))=
 \begin{cases}
 u-\lfloor u\rfloor-\frac12,&u\notin\ZZ,\\
 0,&u\in\ZZ.
 \end{cases}
\end{equation}
The Dedekind sum is
\begin{equation}\label{eq:dedekind-sum}
 s(h,k)=\sum_{j=1}^{k-1}\left(\!\left(\frac jk\right)\!\right)
 \left(\!\left(\frac{hj}{k}\right)\!\right).
\end{equation}
We use the real normalization
\begin{equation}\label{eq:Ak-def}
 A_k(n)=\sum_{\substack{1\le h\le k\\(h,k)=1}}
 \cos\!\left(\pi\left(s(h,k)-\frac{2nh}{k}\right)\right).
\end{equation}
The equivalent complex-exponential version may be used throughout.  Two elementary properties will repeatedly matter:
\begin{equation}\label{eq:Ak-properties}
 A_k(n+k)=A_k(n),\qquad |A_k(n)|\le\varphi(k)\le k.
\end{equation}
The first follows because shifting $n$ by $k$ changes every cosine argument by an integral multiple of $2\pi$; the second is the triangle inequality.  The first few amplitudes are
\begin{equation}\label{eq:Ak-first}
 A_1(n)=1,\qquad A_2(n)=(-1)^n,
\end{equation}
\begin{equation}
 A_3(n)=2\cos\!\left(\frac{2\pi n}{3}-\frac{\pi}{18}\right),\qquad
 A_4(n)=2\cos\!\left(\frac{\pi n}{2}-\frac{\pi}{8}\right).
\end{equation}
The periodicity in \cref{eq:Ak-properties} is the source of the arithmetic phase in the inverse transseries.

\subsection{Bell and Stirling conventions}

The exponential partial Bell polynomials $B_{m,q}$ are defined by
\begin{equation}\label{eq:partial-bell-def}
 \frac1{q!}\left(\sum_{j\ge1}x_j\frac{z^j}{j!}\right)^q
 =\sum_{m\ge q}B_{m,q}(x_1,\ldots,x_{m-q+1})\frac{z^m}{m!}.
\end{equation}
The complete Bell polynomial is
\begin{equation}\label{eq:complete-bell-def}
 Y_m(x_1,\ldots,x_m)=\sum_{q=0}^{m}B_{m,q}(x_1,\ldots,x_{m-q+1}).
\end{equation}
Consequently, if $F(z)=\exp(\sum_{j\ge1}\ell_jz^j)$, then
\begin{equation}\label{eq:bell-exponential}
 [z^m]F(z)=\frac1{m!}Y_m(1!\ell_1,2!\ell_2,\ldots,m!\ell_m).
\end{equation}
This is precisely the exponential formula used in combinatorial coefficient calculus.

The signed Stirling numbers of the first kind $s(m,q)$ are fixed by
\begin{equation}\label{eq:stirling-signed}
 u(u-1)\cdots(u-m+1)=\sum_{q=0}^{m}s(m,q)u^q,
\end{equation}
and equivalently
\begin{equation}\label{eq:log-stirling}
 \bigl(\log(1+z)\bigr)^q=q!\sum_{m\ge q}s(m,q)\frac{z^m}{m!}.
\end{equation}

\subsection{What ``full transseries'' means here}

A single one-variable Hahn series is not the most faithful global model for the Rademacher expansion.  The relative positive-sector exponents
\[
 1-\frac1k\quad(k\ge2)
\]
increase to $1$, while the reflected exponents
\[
 1+\frac1k\quad(k\ge2)
\]
decrease to $1$.  In particular, the latter set is not well ordered in the usual asymptotic ordering.  We therefore use the following two complementary notions.

\begin{definition}[Exact analytic transseries]
The full forward transseries is the convergent Rademacher sum \cref{eq:intro-exact}, regarded as an exact sum of exponential sectors.
\end{definition}

\begin{definition}[Finite-level formal transseries]
For $K\ge1$, retain the sectors with $1\le k\le K$ and treat their distinct exponentials as independent formal variables.  All coefficient extraction and inversion is then performed in an ordinary finitely generated multivariate power-series ring.  The limit $K\to\infty$ is taken analytically through Rademacher convergence, not by pretending that the entire support is a locally finite one-variable Hahn grid.
\end{definition}

This finite-level viewpoint is more than a technical convenience: it is what makes the periodic amplitudes $A_k(n)$ compatible with inversion.

\section{Rademacher's formula in exact transseries normal form}\label{sec:exact-forward}

We take as input Rademacher's convergent formula
\begin{equation}\label{eq:rademacher-input}
 p(n)=\frac1{\pi\sqrt2}\sum_{k\ge1}A_k(n)\sqrt{k}\,
 \frac{\dd}{\dd n}
 \left(\frac{\sinh(x/k)}{\sqrt N}\right),
\end{equation}
with the variables in \cref{eq:master-variables}.  This normalization agrees with the modern formalization in~\cite{EberlPaulson2026}.

\begin{theorem}[Exact elementary sector decomposition]\label{thm:exact-sector}
For every integer $n\ge1$,
\begin{equation}\label{eq:exact-sector}
 p(n)=\frac{\Ccal}{x^2}\sum_{k=1}^{\infty}\frac{A_k(n)}{\sqrt{k}}
 \left\{
 \e^{x/k}\left(1-\frac{k}{x}\right)
 +\e^{-x/k}\left(1+\frac{k}{x}\right)
 \right\}.
\end{equation}
The series is convergent and has exactly the same sum as Rademacher's series.
\end{theorem}

\begin{proof}
Since $\dd x/\dd n=x/(2N)$,
\begin{align}
 \frac{\dd}{\dd n}\left(N^{-1/2}\sinh(x/k)\right)
 &=\frac1{2N^{3/2}}
 \left(\frac{x}{k}\cosh(x/k)-\sinh(x/k)\right)\\
 &=\frac{x}{4kN^{3/2}}
 \left\{
 \e^{x/k}\left(1-\frac{k}{x}\right)
 +\e^{-x/k}\left(1+\frac{k}{x}\right)
 \right\}.
\end{align}
Multiplying by $A_k(n)\sqrt{k}/(\pi\sqrt2)$ gives
\[
 \frac{A_k(n)}{4\sqrt{3k}\,N}
 \left\{
 \e^{x/k}\left(1-\frac{k}{x}\right)
 +\e^{-x/k}\left(1+\frac{k}{x}\right)
 \right\},
\]
because $x/(\pi\sqrt2\sqrt N)=1/\sqrt3$.  Finally,
\[
 \frac1{4\sqrt{3k}\,N}=\frac{\Ccal}{\sqrt{k}\,x^2}
\]
by $x^2=(2\pi^2/3)N$.  The transformation is termwise algebra, so convergence is inherited from \cref{eq:rademacher-input}.
\end{proof}

\begin{corollary}[Dominant sector and exact relative correction]\label{cor:relative-correction}
Define
\begin{equation}\label{eq:Pplus}
 P_+(x)=\Ccal\,\e^x x^{-2}\left(1-\frac1x\right).
\end{equation}
Then
\begin{equation}\label{eq:p-relative}
 p(n)=P_+(x)\bigl(1+\Ecal(x;n)\bigr),
\end{equation}
where
\begin{align}
 \Ecal(x;n)
 &=\e^{-2x}\frac{1+x^{-1}}{1-x^{-1}}\label{eq:E-full}\\
 &\quad+\sum_{k\ge2}\frac{A_k(n)}{\sqrt{k}}
 \left\{
 \e^{-(1-1/k)x}\frac{1-k/x}{1-1/x}
 +\e^{-(1+1/k)x}\frac{1+k/x}{1-1/x}
 \right\}.
 \nonumber
\end{align}
\end{corollary}

\begin{proof}
Separate the $k=1$ positive exponential in \cref{eq:exact-sector} and divide all remaining terms by it.
\end{proof}

\begin{resultbox}[Interpretation]
The leading Hardy--Ramanujan approximation is the first approximation to $P_+(x)$ after replacing $N$ by $n$ and discarding $1-1/x$.  In contrast, \cref{eq:p-relative,eq:E-full} records every Rademacher arc, both growing and reflected, with its exact algebraic amplitude.  It is therefore a complete exponentially improved expansion rather than merely an all-orders Poincare series.
\end{resultbox}

\section{Ordering of the sectors and a direct tail estimate}\label{sec:tail}

For a fixed $k$, the growing sector has absolute exponential scale $\e^{x/k}$ and the reflected sector has scale $\e^{-x/k}$.  Relative to $P_+(x)$, their exponential gaps are
\begin{equation}\label{eq:gaps}
 \delta_{k,+}=1-\frac1k\quad(k\ge2),\qquad
 \delta_{k,-}=1+\frac1k\quad(k\ge2),\qquad
 \delta_{1,-}=2.
\end{equation}
The first correction is therefore the $k=2$ growing sector, of relative order $\e^{-x/2}$, with parity amplitude $(-1)^n/\sqrt2$.

The following elementary estimate is less sharp than Lehmer's classical remainder bounds but is sufficient to justify the hierarchy used in this article.

\begin{theorem}[Fixed-level Rademacher tail]\label{thm:rough-tail}
Let $K\ge1$ be fixed and let $\Rcal_K(n)$ be the part of \cref{eq:exact-sector} with $k>K$.  As $x\to\infty$,
\begin{equation}\label{eq:rough-tail-absolute}
 \Rcal_K(n)=\bigO_K\!\left(x^{-1/2}\e^{x/(K+1)}\right),
\end{equation}
and therefore
\begin{equation}\label{eq:rough-tail-relative}
 \frac{\Rcal_K(n)}{P_+(x)}
 =\bigO_K\!\left(x^{3/2}
 \e^{-(1-1/(K+1))x}\right).
\end{equation}
\end{theorem}

\begin{proof}
Use $|A_k(n)|\le k$.  For $K<k\le x$, we have $k/x\le1$, hence the brace in \cref{eq:exact-sector} is bounded in absolute value by a constant multiple of $\e^{x/k}\le\e^{x/(K+1)}$.  The sum of the prefactors is
\[
 x^{-2}\sum_{k\le x}\sqrt{k}=\bigO(x^{-1/2}),
\]
which proves the desired bound for this range.

For $k>x$, put $z=x/k\in(0,1)$.  The brace is
\[
 2\left(\cosh z-\frac{\sinh z}{z}\right)=\bigO(z^2)
\]
uniformly for $0<z\le1$.  Its full $k$th contribution is therefore
\[
 \bigO\!\left(x^{-2}\sqrt{k}\frac{x^2}{k^2}\right)
 =\bigO(k^{-3/2}).
\]
Summing over $k>x$ gives $\bigO(x^{-1/2})$, which is absorbed by \cref{eq:rough-tail-absolute}.  Dividing by $P_+(x)\asymp x^{-2}\e^x$ proves \cref{eq:rough-tail-relative}.
\end{proof}

\begin{remark}
For a specified residue class, some amplitudes $A_k(n)$ may vanish, so the first nonzero omitted sector can be smaller than the generic estimate.  Lehmer's bounds and subsequent refinements provide much sharper constants.  The role of \cref{thm:rough-tail} is structural: every fixed Rademacher level is exponentially accurate relative to the dominant sector.
\end{remark}

\section{All-order coefficients in the unshifted variable}\label{sec:forward-coefficients}

The shifted variable $x$ makes each sector elementary, but many applications prefer the conventional scale $n^{-1/2}$.  The infinite algebraic coefficient arrays arise entirely from re-expanding the shift $N=n-1/24$.

Put
\begin{equation}\label{eq:t-def}
 t=n^{-1/2},\qquad a=\frac1{24},\qquad
 N=n(1-at^2),\qquad
 x=\frac{\lambda}{t}\sqrt{1-at^2}.
\end{equation}
For $\sigma\in\{+1,-1\}$, define
\begin{equation}\label{eq:Cksigma}
 C_{k,\sigma}(t)=
 (1-at^2)^{-1}
 \exp\!\left[
 \frac{\sigma\lambda}{kt}
 \left(\sqrt{1-at^2}-1\right)
 \right]
 \left(1-\frac{\sigma kt}{\lambda\sqrt{1-at^2}}\right).
\end{equation}
The apparent singularity at $t=0$ is removable because
$\sqrt{1-at^2}-1=\bigO(t^2)$.

\begin{theorem}[Universal sector coefficient generator]\label{thm:universal-sector-gf}
For every fixed $k\ge1$ and $\sigma\in\{+1,-1\}$, let
\begin{equation}\label{eq:c-def}
 C_{k,\sigma}(t)=\sum_{m\ge0}c_{k,\sigma,m}t^m,
 \qquad c_{k,\sigma,0}=1.
\end{equation}
Then the corresponding term in the exact Rademacher expansion is
\begin{equation}\label{eq:sector-unshifted}
 P_{k,\sigma}(n)
 =\frac{A_k(n)}{4\sqrt{3k}\,n}\,
 \e^{\sigma\lambda\sqrt n/k}
 \sum_{m\ge0}c_{k,\sigma,m}n^{-m/2}.
\end{equation}
The series in \cref{eq:c-def} is the Taylor series of \cref{eq:Cksigma}; in particular it converges for all sufficiently small $|t|$ and supplies an asymptotic expansion to arbitrary order.
\end{theorem}

\begin{proof}
The $(k,\sigma)$ term in \cref{eq:exact-sector} is
\[
 \frac{A_k(n)}{4\sqrt{3k}\,N}\,
 \e^{\sigma x/k}\left(1-\frac{\sigma k}{x}\right).
\]
Insert \cref{eq:t-def}, factor out
$A_k(n)\e^{\sigma\lambda/tk}/(4\sqrt{3k}\,n)$, and collect the remaining factors.  The result is exactly \cref{eq:Cksigma}.  Analyticity at $t=0$ follows from the removable singularity just noted and the binomial series for $\sqrt{1-at^2}$.
\end{proof}

\subsection{Logarithmic coefficients and a Bell-polynomial formula}

Write
\begin{equation}\label{eq:ell-def}
 \log C_{k,\sigma}(t)=\sum_{j\ge1}\ell_{k,\sigma,j}t^j.
\end{equation}
All $\ell_{k,\sigma,j}$ can be written as finite sums.

\begin{theorem}[Explicit logarithmic coefficients]\label{thm:ell-explicit}
For $j\ge1$,
\begin{align}
 \ell_{k,\sigma,j}
 &=\mathbf 1_{2\mid j}\frac{a^{j/2}}{j/2}
 \label{eq:ell-explicit}\\
 &\quad+\mathbf 1_{j\text{ odd}}
 \frac{\sigma\lambda}{k}
 \binom{1/2}{(j+1)/2}(-a)^{(j+1)/2}
 \nonumber\\
 &\quad-
 \sum_{\substack{1\le q\le j\\j-q\text{ even}}}
 \frac1q\left(\frac{\sigma k}{\lambda}\right)^q
 \frac{\rising{q/2}{(j-q)/2}}{((j-q)/2)!}
 a^{(j-q)/2}.
 \nonumber
\end{align}
Consequently,
\begin{equation}\label{eq:c-bell}
 c_{k,\sigma,m}
 =\frac1{m!}Y_m\!\left(
 1!\ell_{k,\sigma,1},2!\ell_{k,\sigma,2},\ldots,
 m!\ell_{k,\sigma,m}\right).
\end{equation}
\end{theorem}

\begin{proof}
Taking the logarithm of \cref{eq:Cksigma} gives
\begin{align*}
 \log C_{k,\sigma}(t)
 &=-\log(1-at^2)
 +\frac{\sigma\lambda}{kt}\bigl(\sqrt{1-at^2}-1\bigr)\\
 &\quad+\log\!\left(1-
 \frac{\sigma kt}{\lambda}(1-at^2)^{-1/2}\right).
\end{align*}
The first term contributes $a^r/r$ at degree $2r$.  The second contributes
\[
 \frac{\sigma\lambda}{k}\binom{1/2}{r}(-a)^r
\]
at degree $2r-1$.  For the third term, use
\begin{align*}
 \log(1-u)&=-\sum_{q\ge1}\frac{u^q}{q},\\
 (1-at^2)^{-q/2}
 &=\sum_{r\ge0}\frac{\rising{q/2}{r}}{r!}a^rt^{2r}.
\end{align*}
Setting $j=q+2r$ gives \cref{eq:ell-explicit}.  Formula \cref{eq:c-bell} is the exponential formula \cref{eq:bell-exponential}.
\end{proof}

\begin{corollary}[Triangular coefficient recurrence]\label{cor:c-recurrence}
The same coefficients satisfy
\begin{equation}\label{eq:c-recurrence}
 c_{k,\sigma,0}=1,\qquad
 c_{k,\sigma,m}=\frac1m\sum_{j=1}^{m}
 j\ell_{k,\sigma,j}\,c_{k,\sigma,m-j}
 \quad(m\ge1).
\end{equation}
\end{corollary}

\begin{proof}
Differentiate $C_{k,\sigma}=\exp(\sum_{j\ge1}\ell_{k,\sigma,j}t^j)$ and compare coefficients in
$C'=(\log C)'C$.
\end{proof}

\begin{remark}[Coefficient-field structure]
Every $c_{k,\sigma,m}$ lies in
$\QQ[\lambda,\lambda^{-1},k,k^{-1}]$ after imposing $\sigma^2=1$ and $a=1/24$.  More precisely, the parity of $m$ constrains the powers of $\lambda$: even $m$ gives a polynomial in $\lambda^2$ and $\lambda^{-2}$, while odd $m$ has an overall factor compatible with $\sigma\lambda^{2r+1}$ or $\sigma\lambda^{-(2r+1)}$.  This follows directly from \cref{eq:ell-explicit} and the Bell formula.
\end{remark}

\subsection{The dominant Poincare expansion}

Only the $(k,\sigma)=(1,+1)$ sector survives at every algebraic order relative to the leading exponential.  Put
\begin{equation}\label{eq:omega-def}
 \omega_m=c_{1,+,m}.
\end{equation}
Then
\begin{equation}\label{eq:dominant-poincare}
 p(n)\sim \frac{\e^{\lambda\sqrt n}}{4\sqrt3\,n}
 \sum_{m\ge0}\omega_m n^{-m/2}.
\end{equation}
For each fixed $M$ the difference after truncation at $m=M-1$ is the sum of an algebraic remainder of order $n^{-M/2}$ in the dominant sector and an exponentially small Rademacher remainder.  The all-order statement with explicit error estimates is studied in~\cite{OSullivan2023,BanerjeeEtAl2022}.

The first coefficients are
\begin{align}
 \omega_0&=1,\nonumber\\
 \omega_1&=-\frac{\sqrt6}{2\pi}-\frac{\sqrt6\,\pi}{144},\nonumber\\
 \omega_2&=\frac1{16}+\frac{\pi^2}{6912},\nonumber\\
 \omega_3&=-\frac{\sqrt6}{32\pi}
 -\frac{\sqrt6\,\pi}{2304}
 -\frac{\sqrt6\,\pi^3}{2985984},\label{eq:omega-first}\\
 \omega_4&=\frac5{1536}+\frac{5\pi^2}{497664}
 +\frac{\pi^4}{286654464}.\nonumber
\end{align}
The formulas above are useful for checking implementations, but the finite closed sum in the next section is better for direct evaluation.

\section{A finite closed formula for the dominant coefficients}\label{sec:osullivan}

O'Sullivan's coefficient formula can be derived directly from the generating function \cref{eq:Cksigma}.  The derivation is included because it reveals a compact Catalan substitution and links the partition expansion to Lagrange inversion.

\begin{theorem}[Finite formula for $\omega_m$]\label{thm:omega-closed}
For every $m\ge0$,
\begin{equation}\label{eq:omega-closed}
 \boxed{
 \omega_m=\frac1{(-4\sqrt6)^m}
 \sum_{j=0}^{\lfloor(m+1)/2\rfloor}
 \binom{m+1}{j}
 \frac{m+1-j}{(m+1-2j)!}
 \left(\frac{\pi}{6}\right)^{m-2j}.}
\end{equation}
Terms with a negative factorial index are understood to be absent.  Formula \cref{eq:omega-closed} is equivalent to the coefficient expression in~\cite{OSullivan2023,BanerjeeEtAl2022}.
\end{theorem}

\begin{proof}
Let
\begin{equation}
 r=\frac{t}{4\sqrt6},\qquad A=\frac{\pi}{6},\qquad
 s=\sqrt{1-4r^2},\qquad
 w=\frac{1-s}{2r}.
\end{equation}
Then
\begin{equation}\label{eq:Catalan-w}
 w=r(1+w^2),\qquad
 r=\frac{w}{1+w^2},\qquad
 s=\frac{1-w^2}{1+w^2}.
\end{equation}
A direct substitution into $C_{1,+}(t)$ gives
\begin{equation}\label{eq:C-Fw}
 C_{1,+}(4\sqrt6\,r)
 =\e^{-Aw}\frac{(1+w^2)^2}{(1-w^2)^2}
 \left(1-\frac{2w}{A(1-w^2)}\right).
\end{equation}
Since
\[
 \frac{\dd w}{\dd r}=\frac{(1+w^2)^2}{1-w^2},
\]
the right-hand side of \cref{eq:C-Fw} is
\begin{equation}\label{eq:C-derivative-H}
 -\frac1A\frac{\dd}{\dd r}
 \left(\frac{\e^{-Aw}}{1-w^2}\right).
\end{equation}
Coefficient extraction therefore yields
\begin{equation}\label{eq:omega-proof-step1}
 [r^m]C_{1,+}(4\sqrt6\,r)
 =-\frac{m+1}{A}[r^{m+1}]
 \frac{\e^{-Aw}}{1-w^2}.
\end{equation}
To extract the remaining coefficient, use $r=w/(1+w^2)$ in Cauchy's coefficient integral.  The Jacobian is
\[
 \dd r=\frac{1-w^2}{(1+w^2)^2}\dd w,
\]
so the factor $1-w^2$ cancels and
\begin{equation}\label{eq:omega-proof-residue}
 [r^{m+1}]\frac{\e^{-Aw}}{1-w^2}
 =[w^{m+1}]\e^{-Aw}(1+w^2)^m.
\end{equation}
Expanding the last expression gives
\begin{align*}
 [r^m]C_{1,+}(4\sqrt6\,r)
 &=-\frac{m+1}{A}
 \sum_{j=0}^{\lfloor(m+1)/2\rfloor}
 \binom{m}{j}
 \frac{(-A)^{m+1-2j}}{(m+1-2j)!}\\
 &=(-1)^m\sum_{j=0}^{\lfloor(m+1)/2\rfloor}
 \binom{m+1}{j}
 \frac{m+1-j}{(m+1-2j)!}A^{m-2j},
\end{align*}
where we used
$(m+1)\binom mj=(m+1-j)\binom{m+1}j$.
Finally,
$[t^m]=(4\sqrt6)^{-m}[r^m]$, proving \cref{eq:omega-closed}.
\end{proof}

\begin{remark}[Why a finite sum exists]
The square-root shift might suggest an infinite convolution at each order.  The Catalan coordinate $w$ absorbs that square root, and the Rademacher derivative creates the exact derivative in \cref{eq:C-derivative-H}.  After the coefficient-change Jacobian cancels, only the polynomial $(1+w^2)^m$ remains.  This is the structural reason the coefficient at a fixed order is a finite binomial sum.
\end{remark}

\section{The multiplicative reciprocal \texorpdfstring{$1/p(n)$}{1/p(n)}}\label{sec:reciprocal}

Although the main inverse in this article is functional, the reciprocal transseries is an immediate corollary of the exact factorization \cref{eq:p-relative}.

\begin{theorem}[Exact reciprocal transseries]\label{thm:reciprocal}
For $x$ sufficiently large that $|\Ecal(x;n)|<1$,
\begin{equation}\label{eq:reciprocal-exact}
 \frac1{p(n)}
 =\frac{x^2}{\Ccal}\e^{-x}
 \frac1{1-x^{-1}}
 \sum_{r\ge0}(-1)^r\Ecal(x;n)^r.
\end{equation}
This identity is convergent.  As a formal transseries, it remains valid without the analytic inequality.
\end{theorem}

\begin{proof}
Invert \cref{eq:p-relative} and expand $(1+\Ecal)^{-1}$ geometrically.
\end{proof}

At a finite Rademacher level, write
\begin{equation}\label{eq:E-linear-q}
 \Ecal_K(x;n)=\sum_{j\in J_K}g_j(x;n)q_j,
\end{equation}
where the $q_j$ are independent exponential variables.  For a multi-index
$\alpha=(\alpha_j)_{j\in J_K}$, put
$|\alpha|=\sum_j\alpha_j$, $\alpha!=\prod_j\alpha_j!$, and
$q^\alpha=\prod_jq_j^{\alpha_j}$.

\begin{corollary}[Closed multinomial coefficient formula]\label{cor:reciprocal-multi}
At finite level $K$,
\begin{equation}\label{eq:reciprocal-multi}
 [q^\alpha](1+\Ecal_K)^{-1}
 =(-1)^{|\alpha|}\frac{|\alpha|!}{\alpha!}
 \prod_{j\in J_K}g_j(x;n)^{\alpha_j}.
\end{equation}
\end{corollary}

\begin{proof}
Only the term $(-1)^{|\alpha|}\Ecal_K^{|\alpha|}$ in the geometric series can produce $q^\alpha$, and the multinomial theorem gives the stated coefficient.
\end{proof}

For the ordinary Poincare expansion, define
\begin{equation}
 \left(\sum_{m\ge0}\omega_mt^m\right)^{-1}
 =\sum_{m\ge0}\eta_mt^m.
\end{equation}
Then
\begin{equation}\label{eq:eta-recurrence}
 \eta_0=1,\qquad
 \eta_m=-\sum_{j=1}^{m}\omega_j\eta_{m-j},
\end{equation}
and
\begin{equation}\label{eq:reciprocal-poincare}
 \frac1{p(n)}\sim 4\sqrt3\,n\,\e^{-\lambda\sqrt n}
 \sum_{m\ge0}\eta_mn^{-m/2}.
\end{equation}
Equivalently, if $\log(\sum\omega_mt^m)=\sum\ell_{1,+,j}t^j$, then
\begin{equation}
 \eta_m=\frac1{m!}Y_m(-1!\ell_{1,+,1},\ldots,-m!\ell_{1,+,m}).
\end{equation}
Thus reciprocal coefficients require no new analytic input.

\section{What is the inverse of a discrete sequence?}\label{sec:meaning-inverse}

The values $p(0)=p(1)=1$ coincide, but $p(n)$ is strictly increasing for $n\ge1$.  Indeed, appending a part $1$ injects the partitions of $n$ into those of $n+1$, and the one-part partition $(n+1)$ is not in the image.  Hence the generalized inverse
\begin{equation}\label{eq:generalized-inverse}
 \nu(y)=\min\{n\ge1:p(n)\ge y\},\qquad y\ge1,
\end{equation}
is well defined and integer valued.

There are three related objects, which should not be conflated.

\begin{enumerate}
\item \emph{The exact generalized inverse $\nu$.}  This is a step function.  It has unit jumps and therefore requires an explicit rounding or sawtooth layer; no purely smooth transseries can represent those jumps uniformly to $\smallO(1)$.
\item \emph{The inverse along the range.}  If $y=p(n)$, the desired answer is the integer $n$.  A smooth asymptotic expression can approximate it with exponentially small error, after which nearest-integer rounding is exact for all sufficiently large $n$.
\item \emph{The inverse of an interpolation.}  Any strictly increasing continuous interpolation $\Pcal$ with $\Pcal(n)=p(n)$ has an ordinary inverse $\Pcal^{-1}$.  Then
\begin{equation}\label{eq:ceil-interpolation}
 \nu(y)=\left\lceil\Pcal^{-1}(y)\right\rceil.
\end{equation}
Different interpolations differ within the final unit-scale layer, so their smooth inverses are not canonically identical.
\end{enumerate}

The Rademacher coefficients create a second issue: $A_k(n)$ is periodic rather than slowly varying.  A naive replacement of $n$ by a real variable inside the trigonometric formula gives one possible analytic interpolation, but not a unique one.  Our solution is intrinsic at each exponential level: fix the residue of $n$ modulo the periods of the retained sectors.  The resulting finite-level function is analytic in the phase variable $x$ and has a canonical asymptotic inverse.  This construction is developed in \cref{sec:full-inverse}.

\begin{warningbox}[Scope of the inverse formulas]
The formulas through \cref{sec:dominant-inverse} are phase independent and therefore apply directly to large real ordinates $y$.  The exponentially small formulas in \cref{sec:full-inverse,sec:first-inverse-sectors} describe the inverse along a specified arithmetic phase, in particular along $y=p(n)$ after the residue class of $n$ is known.  The phase can itself be recovered by rounding the phase-independent dominant inverse for all sufficiently large exact ordinates; see \cref{thm:eventual-rounding}.
\end{warningbox}

\section{Exact Lambert carrier and its polynomial--logarithmic expansion}\label{sec:Lambert-core}

The first inversion step retains only
\begin{equation}\label{eq:carrier}
 y=\Ccal\frac{\e^X}{X^2}.
\end{equation}
This differs from the complete dominant sector $P_+(x)$ by the factor $1-1/x$, which will be restored in \cref{sec:dominant-inverse}.

\begin{theorem}[Exact Lambert carrier]\label{thm:Lambert-carrier}
For sufficiently large $y$, the large positive solution of \cref{eq:carrier} is
\begin{equation}\label{eq:X-Lambert}
 \boxed{
 X(y)=-2W_{-1}\!\left(-\frac12\sqrt{\frac{\Ccal}{y}}\right).}
\end{equation}
Equivalently, with
\begin{equation}\label{eq:R-L}
 R=\log\frac{y}{\Ccal},\qquad L=\log R,
\end{equation}
we have
\begin{equation}\label{eq:X-phase}
 X-2\log X=R.
\end{equation}
\end{theorem}

\begin{proof}
Equation \cref{eq:carrier} implies
\[
 \sqrt{\frac{\Ccal}{y}}=X\e^{-X/2}.
\]
Multiplying by $-1/2$ gives
\[
 -\frac12\sqrt{\frac{\Ccal}{y}}
 =\left(-\frac X2\right)\e^{-X/2}.
\]
The solution with $X\to+\infty$ corresponds to the real branch $W_{-1}$, proving \cref{eq:X-Lambert}.  Taking logarithms of \cref{eq:carrier} proves \cref{eq:X-phase}.
\end{proof}

The Lambert formula is already an exact closed form for the carrier.  To obtain a pure logarithmic transseries, write $X=R(1+v)$.  Equation \cref{eq:X-phase} becomes
\begin{equation}\label{eq:v-Lagrange}
 v=z\bigl(L+\log(1+v)\bigr),\qquad z=\frac2R.
\end{equation}
This is in the standard Lagrange form $v=z\phi(v)$.

\begin{theorem}[All polynomial--logarithmic carrier coefficients]\label{thm:X-polynomials}
As $y\to\infty$,
\begin{equation}\label{eq:X-full-log}
 X\sim R+2L+\sum_{j\ge1}\frac{P_j(L)}{R^j},
\end{equation}
where
\begin{align}
 P_j(L)
 &=\frac{2^{j+1}}{j+1}[u^j]
 \bigl(L+\log(1+u)\bigr)^{j+1}
 \label{eq:Pj-coeff}\\
 &=2^{j+1}\sum_{q=1}^{j}
 \frac{s(j,q)}{(j+1-q)!}L^{j+1-q}.
 \label{eq:Pj-stirling}
\end{align}
For every fixed $J$,
\begin{equation}\label{eq:X-remainder}
 X=R+2L+\sum_{j=1}^{J}\frac{P_j(L)}{R^j}
 +\bigO\!\left(\frac{L^{J+1}}{R^{J+1}}\right).
\end{equation}
\end{theorem}

\begin{proof}
Lagrange--Buer\-mann inversion applied to \cref{eq:v-Lagrange} gives
\begin{equation}
 v=\sum_{m\ge1}\frac{z^m}{m}[u^{m-1}]
 \bigl(L+\log(1+u)\bigr)^m.
\end{equation}
Multiplying by $R$ and putting $m=j+1$ yields \cref{eq:Pj-coeff}.  Expand the power by the binomial theorem and use \cref{eq:log-stirling}.  Only $1\le q\le j$ contributes to $[u^j]$, and simplification gives \cref{eq:Pj-stirling}.  The remainder estimate follows from the analytic implicit-function theorem applied to \cref{eq:v-Lagrange}, with $z=2/R$ and $L=\log R$ treated on the asymptotic scale $L/R\to0$.
\end{proof}

The first polynomials are
\begin{align}
 P_1(L)&=4L,\nonumber\\
 P_2(L)&=-4L(L-2),\nonumber\\
 P_3(L)&=\frac83L(2L^2-9L+6),\label{eq:P-first}\\
 P_4(L)&=-\frac83L(3L^3-22L^2+36L-12),\nonumber\\
 P_5(L)&=\frac{16}{15}L(12L^4-125L^3+350L^2-300L+60).\nonumber
\end{align}
Thus, for example,
\begin{align}\label{eq:X-explicit-five}
 X={}&R+2L+\frac{4L}{R}
 +\frac{-4L^2+8L}{R^2}
 +\frac{\frac{16}{3}L^3-24L^2+16L}{R^3}\\
 &+\frac{-8L^4+\frac{176}{3}L^3-96L^2+32L}{R^4}
 +\cdots.\nonumber
\end{align}

\begin{remark}[Two useful normalizations]
One may instead expand $-W_{-1}(-\varepsilon)$ in
$H=\log(1/\varepsilon)$ and $\log H$.  Formula \cref{eq:Pj-stirling} is tailored to the partition problem because $R=\log(y/\Ccal)$ makes the phase equation exactly $X-2\log X=R$ and avoids additional constant translations.
\end{remark}

\section{Inverting the full dominant Rademacher sector}\label{sec:dominant-inverse}

We now solve
\begin{equation}\label{eq:dominant-inverse-equation}
 y=P_+(x)=\Ccal\e^x x^{-2}\left(1-\frac1x\right).
\end{equation}
The Lambert carrier $X$ solves the same equation without the final factor.  The correction is algebraic in $1/X$.

\begin{theorem}[Rational correction series]\label{thm:U-series}
Let $X=X(y)$ be given by \cref{eq:X-Lambert}, put $z=X^{-1}$, and let $U(z)$ be the unique formal series in $z\QQ[[z]]$ satisfying
\begin{equation}\label{eq:U-implicit}
 \boxed{U+\log(1-z+zU)-3\log(1+zU)=0.}
\end{equation}
Then the asymptotic inverse of \cref{eq:dominant-inverse-equation} is
\begin{equation}\label{eq:x0-XU}
 x_0(y)=X+U(X^{-1}).
\end{equation}
Moreover, $U$ is analytic in a neighborhood of $z=0$, so its series has a positive radius of convergence.
\end{theorem}

\begin{proof}
Put $x=X+U$ and $z=1/X$.  Dividing \cref{eq:dominant-inverse-equation} by the carrier identity $y=\Ccal\e^XX^{-2}$ gives
\[
 \e^U(1+zU)^{-2}
 \left(1-\frac{z}{1+zU}\right)=1.
\]
Since
\[
 1-\frac{z}{1+zU}=\frac{1-z+zU}{1+zU},
\]
taking logarithms gives \cref{eq:U-implicit}.  At $(z,U)=(0,0)$ the left-hand side vanishes and its derivative with respect to $U$ equals $1$.  The analytic implicit-function theorem therefore gives a unique analytic solution $U(z)$ near zero.  Substitution recovers \cref{eq:dominant-inverse-equation}.
\end{proof}

\begin{theorem}[Triangular formula for the rational coefficients]\label{thm:b-recurrence}
Write
\begin{equation}\label{eq:U-b}
 U(z)=\sum_{m\ge1}b_mz^m,
 \qquad U_{<m}(z)=\sum_{j=1}^{m-1}b_jz^j.
\end{equation}
Then
\begin{equation}\label{eq:b-recurrence-compact}
 \boxed{
 b_m=-[z^m]\left\{
 \log(1-z+zU_{<m})-3\log(1+zU_{<m})
 \right\}.}
\end{equation}
Equivalently,
\begin{align}
 b_m
 &=\sum_{q=1}^{m}\frac{(-1)^q}{q}
 [z^{m-q}](U_{<m}-1)^q
 \label{eq:b-recurrence-expanded}\\
 &\quad+3\sum_{q=1}^{m}\frac{(-1)^{q+1}}{q}
 [z^{m-q}]U_{<m}^{q}.
 \nonumber
\end{align}
In particular, every $b_m$ is rational.
\end{theorem}

\begin{proof}
In \cref{eq:U-implicit}, the coefficient of $z^m$ contributed by $U$ is $b_m$.  Both logarithmic terms contain $U$ only after multiplication by $z$, so their $z^m$ coefficients involve only $b_1,\ldots,b_{m-1}$.  Solving the coefficient equation gives \cref{eq:b-recurrence-compact}.  Expanding the two logarithms gives \cref{eq:b-recurrence-expanded}.  Rationality follows inductively.
\end{proof}

The first fifteen coefficients are
\begin{align}
 b_1&=1,&
 b_2&=\frac52,&
 b_3&=\frac{13}{3},&
 b_4&=\frac{53}{12},\nonumber\\
 b_5&=-\frac{14}{5},&
 b_6&=-\frac{763}{30},&
 b_7&=-\frac{2179}{35},&
 b_8&=-\frac{42289}{630},\nonumber\\
 b_9&=\frac{132973}{1260},&
 b_{10}&=\frac{3449711}{5040},&
 b_{11}&=\frac{29813143}{18480},&
 b_{12}&=\frac{496569589}{356400},\nonumber\\
 b_{13}&=-\frac{4334124833}{926640},&
 b_{14}&=-\frac{509723063851}{21621600},&
 b_{15}&=-\frac{298798016297}{5896800}.\label{eq:b-first}
\end{align}
Thus
\begin{equation}\label{eq:x0-first}
 x_0=X+\frac1X+\frac{5}{2X^2}+\frac{13}{3X^3}
 +\frac{53}{12X^4}-\frac{14}{5X^5}+\cdots.
\end{equation}

\subsection{The phase-independent inverse index}

Define
\begin{equation}\label{eq:nplus-def}
 n_+(y)=\frac1{24}+\frac{3x_0(y)^2}{2\pi^2}.
\end{equation}
This is the inverse index supplied by the complete dominant positive Rademacher sector.  It already includes every algebraic correction; the difference from an exact index $n$ with $y=p(n)$ is beyond all algebraic orders.

\begin{corollary}[All algebraic coefficients of the inverse index]\label{cor:nplus-expansion}
Let
\begin{equation}\label{eq:beta-def}
 \beta_j=3b_{j+1}+\frac32\sum_{r=1}^{j-1}b_rb_{j-r}
 \qquad(j\ge1).
\end{equation}
Then
\begin{equation}\label{eq:nplus-expansion}
 \boxed{
 n_+(y)=\frac{3X^2}{2\pi^2}
 +\left(\frac1{24}+\frac3{\pi^2}\right)
 +\frac1{\pi^2}\sum_{j\ge1}\frac{\beta_j}{X^j}.}
\end{equation}
The first values are
\begin{align}
 \beta_1&=\frac{15}{2},&
 \beta_2&=\frac{29}{2},&
 \beta_3&=\frac{83}{4},&
 \beta_4&=\frac{559}{40},\nonumber\\
 \beta_5&=-\frac{611}{20},&
 \beta_6&=-\frac{112459}{840},&
 \beta_7&=-\frac{101329}{420},&
 \beta_8&=-\frac{228677}{3360}.\label{eq:beta-first}
\end{align}
\end{corollary}

\begin{proof}
Insert $x_0=X+\sum_{m\ge1}b_mX^{-m}$ into \cref{eq:nplus-def}.  The term $2XU$ contributes $3b_{j+1}/(\pi^2X^j)$, and $U^2$ contributes
$3(\sum_{r=1}^{j-1}b_rb_{j-r})/(2\pi^2X^j)$.
\end{proof}

\begin{remark}[The constant shift]
The carrier-only index
$1/24+3X^2/(2\pi^2)$ underestimates the exact index by a quantity tending to $3/\pi^2$.  The first correction $U\sim X^{-1}$ contributes exactly $3/\pi^2$ to the inverse index.  This explains the approximately $-0.304$ carrier error visible in the numerical table of \cref{sec:numerics}.
\end{remark}

\section{The full arithmetic inverse transseries}\label{sec:full-inverse}

The remaining correction to $n_+(y)$ is exponentially small but arithmetically oscillatory.  The oscillation must be frozen before one can speak of an analytic inverse.

\subsection{Finite Rademacher levels and residue classes}

For $K\ge1$, put
\begin{equation}\label{eq:LK}
 L_K=\operatorname{lcm}(1,2,\ldots,K).
\end{equation}
Fix $r\in\ZZ/L_K\ZZ$.  Since $A_k(n)$ is periodic modulo $k$, the values $A_k(r)$ are unambiguously defined for all $k\le K$.  Define the finite-level analytic function
\begin{equation}\label{eq:PKr}
 P_{K,r}(x)=\frac{\Ccal}{x^2}
 \sum_{k=1}^{K}\frac{A_k(r)}{\sqrt{k}}
 \left\{
 \e^{x/k}\left(1-\frac{k}{x}\right)
 +\e^{-x/k}\left(1+\frac{k}{x}\right)
 \right\}.
\end{equation}
For large $x$, the $k=1$ positive term dominates, so $P_{K,r}'(x)>0$ and $P_{K,r}$ has a unique inverse branch tending to infinity.  Denote it by $x_{K,r}(y)$.

Relative to $P_+$, write
\begin{equation}\label{eq:EKr}
 P_{K,r}(x)=P_+(x)\bigl(1+\Ecal_{K,r}(x)\bigr),
\end{equation}
where
\begin{align}
 \Ecal_{K,r}(x)
 &=\e^{-2x}\frac{1+x^{-1}}{1-x^{-1}}\label{eq:EKr-explicit}\\
 &\quad+\sum_{k=2}^{K}\frac{A_k(r)}{\sqrt{k}}
 \left\{
 \e^{-(1-1/k)x}\frac{1-k/x}{1-1/x}
 +\e^{-(1+1/k)x}\frac{1+k/x}{1-1/x}
 \right\}.
 \nonumber
\end{align}
Introduce the dominant logarithmic phase
\begin{equation}\label{eq:Phi-def}
 \Phi(x)=x-2\log x+\log\left(1-\frac1x\right),
 \qquad x>1.
\end{equation}
Then $x_0=h(R)$, where $R=\log(y/\Ccal)$ and $h=\Phi^{-1}$, while the full finite-level inverse solves
\begin{equation}\label{eq:phase-full}
 \Phi(x)+\Psi_{K,r}(x)=R,
 \qquad
 \Psi_{K,r}(x)=\log\bigl(1+\Ecal_{K,r}(x)\bigr).
\end{equation}
The slope divisor is
\begin{equation}\label{eq:D0}
 D(x)=\Phi'(x)
 =1-\frac2x+\frac1{x(x-1)}
 =\frac{x^2-3x+3}{x(x-1)}>0.
\end{equation}
More generally,
\begin{equation}\label{eq:Phi-derivatives}
 \Phi^{(m)}(x)=\mathbf1_{m=1}
 +(-1)^{m-1}(m-1)!
 \left((x-1)^{-m}-3x^{-m}\right),\qquad m\ge1.
\end{equation}

\subsection{Independent exponential variables}

Let
\begin{equation}\label{eq:JK}
 J_K=\{(1,-)\}\cup\{(k,+),(k,-):2\le k\le K\}.
\end{equation}
For $j\in J_K$, define the gap $\delta_j$ and amplitude $g_j$ by
\begin{align}
 \delta_{1,-}&=2,&
 g_{1,-}(x)&=\frac{1+x^{-1}}{1-x^{-1}},\label{eq:g1minus}\\
 \delta_{k,+}&=1-\frac1k,&
 g_{k,+}(x;r)&=\frac{A_k(r)}{\sqrt{k}}
 \frac{1-k/x}{1-1/x},\label{eq:gkplus}\\
 \delta_{k,-}&=1+\frac1k,&
 g_{k,-}(x;r)&=\frac{A_k(r)}{\sqrt{k}}
 \frac{1+k/x}{1-1/x}.
 \label{eq:gkminus}
\end{align}
Introduce independent variables $q_j$ and write
\begin{equation}\label{eq:E-formal}
 \Ecal_{K,r}(x;\bm q)=\sum_{j\in J_K}q_jg_j(x;r).
\end{equation}
At the end of the computation one specializes
\begin{equation}\label{eq:q-specialization}
 q_j=\e^{-\delta_jx_0}.
\end{equation}
When $x=x_0+v$, the actual exponential is
$q_j\e^{-\delta_jv}$, so define
\begin{equation}\label{eq:Psi-formal-v}
 \Psi_{K,r}(x_0+v;\bm q)
 =\log\!\left(
 1+\sum_{j\in J_K}q_j
 g_j(x_0+v;r)\e^{-\delta_jv}
 \right).
\end{equation}
This formalization treats resonant exponent sums correctly: distinct products are first kept as distinct multivariate monomials and may be combined only after \cref{eq:q-specialization}.

\subsection{The triangular Newton--Hensel recurrence}

Choose any monomial order on $\NN^{J_K}$ that is compatible with addition, for example graded lexicographic order.  For a multi-index $\alpha$, write
$\bm q^\alpha=\prod_jq_j^{\alpha_j}$.  Seek
\begin{equation}\label{eq:v-multi}
 v(\bm q)=\sum_{\alpha\ne0}V_\alpha(x_0;r)\bm q^\alpha.
\end{equation}
For a fixed $\alpha$, let $v_{<\alpha}$ be the sum of terms preceding $\alpha$ in the chosen order.

\begin{theorem}[Complete multivariate inverse-coefficient recurrence]\label{thm:multivariate-recurrence}
Let
\begin{equation}\label{eq:kappa}
 \kappa_m(x_0)=\frac{\Phi^{(m)}(x_0)}{m!},\qquad
 \kappa_1(x_0)=D(x_0).
\end{equation}
The unique formal solution of \cref{eq:phase-full} with $x=x_0+v$ is given recursively by
\begin{equation}\label{eq:Valpha-recurrence}
 \boxed{
 V_\alpha=-\frac1{D(x_0)}[\bm q^\alpha]
 \left\{
 \sum_{m\ge2}\kappa_m(x_0)v_{<\alpha}^{\,m}
 +\Psi_{K,r}(x_0+v_{<\alpha};\bm q)
 \right\}.}
\end{equation}
At each $\alpha$, the right-hand side is a finite expression involving only previously computed coefficients.  It is a rational combination of $x_0$, $(x_0-1)^{-1}$, the gaps $\delta_j$, and the arithmetic amplitudes $A_k(r)/\sqrt{k}$.
\end{theorem}

\begin{proof}
Taylor expansion of \cref{eq:phase-full} at $x_0$ gives
\begin{equation}\label{eq:normalized-v-equation}
 D(x_0)v+\sum_{m\ge2}\kappa_m(x_0)v^m
 +\Psi_{K,r}(x_0+v;\bm q)=0.
\end{equation}
The only term in which $V_\alpha$ can occur linearly at monomial $\bm q^\alpha$ is $D(x_0)v$.  Every term $v^m$ with $m\ge2$ uses a sum of at least two positive multi-indices and therefore depends only on lower coefficients.  Every term in $\Psi$ contains at least one explicit $q_j$; if $V_\alpha$ appears inside $v$, multiplication by that $q_j$ raises the multi-index beyond $\alpha$.  Hence the coefficient equation at $\alpha$ is triangular and solving it gives \cref{eq:Valpha-recurrence}.  Finiteness follows because there are only finitely many variables and a fixed multi-index has finitely many additive decompositions.
\end{proof}

\begin{corollary}[First coefficient in each exponential direction]\label{cor:first-V}
Let $e_j$ be the unit multi-index of the variable $q_j$.  Then
\begin{equation}\label{eq:first-V}
 V_{e_j}(x_0;r)=-\frac{g_j(x_0;r)}{D(x_0)}.
\end{equation}
After \cref{eq:q-specialization}, the first-order contribution of sector $j$ is
\begin{equation}\label{eq:first-v-actual}
 v_j^{(1)}(y;r)=-\frac{g_j(x_0;r)}{D(x_0)}\e^{-\delta_jx_0}.
\end{equation}
\end{corollary}

\subsection{A closed Lagrange--Buer\-mann formula}

The triangular recurrence is ideal for exact coefficient production.  A second formula packages all orders into one operator identity.

\begin{theorem}[Lagrange--Buer\-mann formula for the full inverse]\label{thm:LB-full-inverse}
Let $h=\Phi^{-1}$, $R=\log(y/\Ccal)$, and
\begin{equation}\label{eq:A-phase}
 A(s)=\Psi_{K,r}(h(s)).
\end{equation}
Then, as a formal transseries in the exponential variables,
\begin{equation}\label{eq:LB-phase}
 S=(I+A)^{-1}(R)
 =R+\sum_{m\ge1}\frac{(-1)^m}{m!}
 \left(\frac{\dd}{\dd R}\right)^{m-1}A(R)^m,
\end{equation}
and
\begin{equation}\label{eq:LB-x}
 \boxed{x_{K,r}(y)=h(S).}
\end{equation}
Equations \cref{eq:LB-phase,eq:LB-x} are an all-orders closed formula for the inverse transseries.
\end{theorem}

\begin{proof}
Put $s=\Phi(x)$.  Then $x=h(s)$ and \cref{eq:phase-full} becomes
\[
 s+A(s)=R.
\]
The classical Lagrange--Buer\-mann reversion formula for the near-identity map $s\mapsto s+A(s)$ gives \cref{eq:LB-phase}.  Composing with $h$ gives \cref{eq:LB-x}.
\end{proof}

To turn \cref{eq:LB-x} into coefficient polynomials, introduce a bookkeeping parameter $\tau$ by replacing $A$ with $\tau A$, and set
\begin{equation}\label{eq:dj-def}
 d_m(R)=\frac{(-1)^m}{m!}
 \left(\frac{\dd}{\dd R}\right)^{m-1}A(R)^m.
\end{equation}
Then $S-R=\sum_{m\ge1}d_m\tau^m$.  Fa\`a di Bruno's formula gives the following explicit Bell-polynomial coefficient.

\begin{corollary}[Coefficient at perturbative order $m$]\label{cor:LB-Bell}
If
\begin{equation}
 x_{K,r}(y;\tau)=x_0+\sum_{m\ge1}v_m\tau^m,
\end{equation}
then
\begin{equation}\label{eq:vm-Bell}
 \boxed{
 v_m=\frac1{m!}\sum_{q=1}^{m}h^{(q)}(R)
 B_{m,q}\bigl(1!d_1,2!d_2,\ldots,(m-q+1)!d_{m-q+1}\bigr).}
\end{equation}
The derivatives of $h$ are generated recursively by
\begin{equation}\label{eq:hprime}
 h'(R)=\frac1{\Phi'(x_0)},
\end{equation}
\begin{equation}\label{eq:h-derivative-recurrence}
 h^{(m)}(R)=-\frac1{\Phi'(x_0)}
 \sum_{q=2}^{m}\Phi^{(q)}(x_0)
 B_{m,q}\bigl(h',h'',\ldots,h^{(m-q+1)}\bigr),
 \qquad m\ge2.
\end{equation}
\end{corollary}

\begin{proof}
Expand $h(R+\sum d_m\tau^m)$ by Taylor's theorem and use the defining generating identity \cref{eq:partial-bell-def} for the coefficient of a power of $\sum d_m\tau^m$.  For \cref{eq:h-derivative-recurrence}, differentiate $\Phi(h(R))=R$ $m$ times and solve the Fa\`a di Bruno identity for the term containing $h^{(m)}$.
\end{proof}

\subsection{Low perturbative orders in differential form}

For checking symbolic implementations, consider the one-parameter equation
\begin{equation}\label{eq:tau-equation}
 \Phi(x)+\tau\Psi(x)=R,
 \qquad x=x_0+\tau v_1+\tau^2v_2+\tau^3v_3+\cdots.
\end{equation}
All functions and derivatives below are evaluated at $x_0$, and $D=\Phi'(x_0)$.  Direct coefficient comparison gives
\begin{align}
 v_1&=-\frac{\Psi}{D},\label{eq:v1-diff}\\
 v_2&=\frac{\Psi\Psi'}{D^2}
 -\frac{\Phi''\Psi^2}{2D^3},\label{eq:v2-diff}\\
 v_3&=-\frac{\Psi(\Psi')^2}{D^3}
 -\frac{\Psi^2\Psi''}{2D^3}
 +\frac{3\Phi''\Psi^2\Psi'}{2D^4}
 +\frac{\Phi'''\Psi^3}{6D^4}
 -\frac{(\Phi'')^2\Psi^3}{2D^5}.
 \label{eq:v3-diff}
\end{align}
These formulas collect many exponential monomials at once; \cref{eq:Valpha-recurrence} is preferable when individual sectors or resonances are required.

\subsection{The profinite arithmetic phase}

The residue systems are compatible because $L_K\mid L_{K+1}$.  A compatible sequence
\begin{equation}\label{eq:profinite-phase}
 r_K\in\ZZ/L_K\ZZ,
 \qquad r_{K+1}\equiv r_K\pmod{L_K},
\end{equation}
defines a profinite integer $\widehat r\in\widehat{\ZZ}$.  Conversely, every $\widehat r$ supplies all amplitudes $A_k(\widehat r)$ by reduction modulo $k$.

\begin{theorem}[Arithmetic-phase organization]\label{thm:profinite}
The family of finite-level inverse transseries
\[
 \{x_{K,r_K}(y)\}_{K\ge1}
\]
is compatible under truncation whenever the residues satisfy \cref{eq:profinite-phase}.  Hence the complete arithmetic inverse is naturally a projective family indexed by $\widehat r\in\widehat{\ZZ}$.  For an actual partition index $n$, the diagonal embedding $n\mapsto\widehat n$ gives its phase.
\end{theorem}

\begin{proof}
Truncating the $(K+1)$st construction removes all variables associated with $k=K+1$.  For every retained $k\le K$, compatibility of the residues gives
$A_k(r_{K+1})=A_k(r_K)$.  Thus the truncated defining equations, and consequently their unique triangular solutions, agree.
\end{proof}

\begin{remark}[Why the profinite parameter is not optional]
At algebraic order, all indices share the same expansion.  At relative order $\e^{-x/2}$, parity appears; at order $\e^{-2x/3}$, the residue modulo $3$ appears; and so on.  No finite real-valued smooth parameter records all these congruence classes, while a profinite integer does so canonically.
\end{remark}

\section{First arithmetic inverse sectors and the master formula}\label{sec:first-inverse-sectors}

At finite level $K$, the first total degree in the exponential variables is especially transparent.  Combining \cref{eq:first-v-actual,eq:g1minus,eq:gkplus,eq:gkminus} gives
\begin{align}
 x_{K,r}(y)
 &=x_0-\frac1{D(x_0)}\Biggl[
 \e^{-2x_0}\frac{1+x_0^{-1}}{1-x_0^{-1}}
 \label{eq:first-all-sectors}\\
 &\quad+\sum_{k=2}^{K}\frac{A_k(r)}{\sqrt{k}}
 \left\{
 \e^{-(1-1/k)x_0}\frac{1-k/x_0}{1-1/x_0}
 +\e^{-(1+1/k)x_0}\frac{1+k/x_0}{1-1/x_0}
 \right\}\Biggr]
 +\bigO_{\bm q}(2),\nonumber
\end{align}
where $\bigO_{\bm q}(2)$ denotes monomials of total degree at least two in the independent $q_j$.

The first two growing arithmetic sectors are
\begin{align}
 v_{2,+}^{(1)}
 &=-\frac{(-1)^r}{\sqrt2\,D(x_0)}
 \frac{1-2/x_0}{1-1/x_0}\e^{-x_0/2},
 \label{eq:k2-inverse}\\
 v_{3,+}^{(1)}
 &=-\frac{2\cos(2\pi r/3-\pi/18)}{\sqrt3\,D(x_0)}
 \frac{1-3/x_0}{1-1/x_0}\e^{-2x_0/3}.
 \label{eq:k3-inverse}
\end{align}
Thus parity enters at scale $\e^{-x_0/2}$ and the residue modulo $3$ at scale $\e^{-2x_0/3}$.  More generally, the residue modulo $k$ first appears in the growing $k$th Rademacher sector at scale $\e^{-(1-1/k)x_0}$.

\begin{remark}[Accumulation near relative exponent one]
There are infinitely many first-order growing sectors with gaps $1-1/k<1$.  Therefore an instruction such as ``retain every term through $\e^{-x}$'' does not select a finite set of Rademacher sectors.  A meaningful truncation must specify a Rademacher level $K$, an arithmetic cutoff, or a combined analytic error criterion.  This is another reason to keep the exact convergent sector sum and finite-level formal calculus side by side.
\end{remark}

\subsection{From phase to index}

Let
\begin{equation}\label{eq:nKr}
 n_{K,r}(y)=\frac1{24}+\frac{3x_{K,r}(y)^2}{2\pi^2}.
\end{equation}
If $v=\sum_{\alpha\ne0}V_\alpha\bm q^\alpha$, then
\begin{equation}\label{eq:n-v}
 n_{K,r}=n_+(y)+\frac{3x_0}{\pi^2}v+\frac{3}{2\pi^2}v^2.
\end{equation}
Hence the coefficient of each exponential monomial is explicit.

\begin{corollary}[Inverse-index coefficient formula]\label{cor:index-coeff}
For $\alpha\ne0$, let $N_\alpha$ be the coefficient of $\bm q^\alpha$ in $n_{K,r}-n_+$.  Then
\begin{equation}\label{eq:Nalpha}
 \boxed{
 N_\alpha=\frac3{\pi^2}
 \left(
 x_0V_\alpha+\frac12
 \sum_{\substack{\beta+\gamma=\alpha\\\beta,\gamma\ne0}}
 V_\beta V_\gamma
 \right).}
\end{equation}
In particular,
\begin{equation}\label{eq:N-first}
 N_{e_j}=-\frac{3x_0}{\pi^2D(x_0)}g_j(x_0;r).
\end{equation}
\end{corollary}

\begin{proof}
Expand the square in \cref{eq:n-v} and compare coefficients.
\end{proof}

Combining all layers produces the main inverse formula.

\begin{theorem}[Master nested inverse transseries]\label{thm:master-inverse}
Fix a finite level $K$ and a residue $r\pmod{L_K}$.  Let
\begin{align*}
 X(y)&=-2W_{-1}\!\left(-\frac12\sqrt{\frac{\Ccal}{y}}\right),\\
 x_0(y)&=X(y)+U(X(y)^{-1}),\\
 v_{K,r}(y)&=\sum_{\alpha\ne0}V_\alpha(x_0;r)
 \exp\!\left(-x_0\sum_{j\in J_K}\alpha_j\delta_j\right),
\end{align*}
where $U$ is determined by \cref{eq:b-recurrence-compact} and the $V_\alpha$ by \cref{eq:Valpha-recurrence}.  Then
\begin{equation}\label{eq:master-inverse}
 \boxed{
 n_{K,r}(y)=\frac1{24}+\frac{3}{2\pi^2}
 \bigl(x_0(y)+v_{K,r}(y)\bigr)^2.}
\end{equation}
Substituting the Stirling-number expansion \cref{eq:X-full-log,eq:Pj-stirling} for $X(y)$ turns \cref{eq:master-inverse} into a transseries expressed solely in
\[
 R^{-1},\quad \log R,\quad
 \exp\!\left[-\delta\left(R+2\log R+\cdots\right)\right],
 \qquad R=\log(y/\Ccal),
\]
with every coefficient generated by finite Bell-polynomial and Cauchy-product operations.
\end{theorem}

\begin{proof}
The first line inverts the elementary carrier, the second restores the full dominant positive sector, and the third is the unique finite-level exponential correction of \cref{thm:multivariate-recurrence}.  Equation \cref{eq:intro-n-from-x} then gives the index.  Closure under substitution follows from the coefficient formulas already proved.
\end{proof}

\section{Validity, truncation error, and discrete recovery}\label{sec:validity}

\subsection{Finite-level inverse error along exact partition values}

Let $n$ be an integer, let $x=\lambda\sqrt{n-1/24}$, and fix the compatible residue $r\equiv n\pmod{L_K}$.  The exact partition value satisfies
\begin{equation}
 p(n)=P_{K,r}(x)+\Rcal_K(n).
\end{equation}
For large $x$, logarithmic derivatives of $P_{K,r}$ are bounded away from zero because the dominant term has slope tending to $1$.

\begin{theorem}[Error after inversion of $K$ Rademacher levels]\label{thm:inverse-tail}
For fixed $K$ and $y=p(n)$,
\begin{equation}\label{eq:xK-error}
 x_{K,r}(y)-x
 =\bigO_K\!\left(
 x^{3/2}\e^{-(1-1/(K+1))x}
 \right),
\end{equation}
and
\begin{equation}\label{eq:nK-error}
 n_{K,r}(y)-n
 =\bigO_K\!\left(
 x^{5/2}\e^{-(1-1/(K+1))x}
 \right).
\end{equation}
\end{theorem}

\begin{proof}
By \cref{thm:rough-tail}, the omitted tail divided by the dominant sector is
$\bigO_K(x^{3/2}\e^{-(1-1/(K+1))x})$.  Taking logarithms changes this by only a comparable amount.  The derivative of $\log P_{K,r}(x)$ tends to $1$, so the mean-value theorem converts the logarithmic residual into the phase error \cref{eq:xK-error}.  Finally,
$n=1/24+3x^2/(2\pi^2)$, whose derivative with respect to $x$ is $3x/\pi^2$, proving \cref{eq:nK-error}.
\end{proof}

The case $K=1$ is particularly useful because it is phase independent.

\begin{corollary}[Beyond-all-orders accuracy of the dominant inverse]\label{cor:dominant-inverse-error}
For $y=p(n)$,
\begin{equation}\label{eq:nplus-error}
 n_+(y)-n=\bigO\!\left(x^{5/2}\e^{-x/2}\right)
 =\bigO\!\left(n^{5/4}\e^{-\lambda\sqrt n/2}\right).
\end{equation}
In particular, $n_+(p(n))-n\to0$ faster than every power of $n^{-1}$.
\end{corollary}

\begin{theorem}[Eventual exact recovery by rounding]\label{thm:eventual-rounding}
There exists $n_0$ such that for every integer $n\ge n_0$,
\begin{equation}\label{eq:rounding-exact}
 n=\operatorname{nint}\bigl(n_+(p(n))\bigr),
\end{equation}
where $\operatorname{nint}$ denotes nearest-integer rounding.  Consequently, for sufficiently large exact ordinates $y=p(n)$, the arithmetic phase $\widehat n$ can be recovered before any exponentially small sector is evaluated.
\end{theorem}

\begin{proof}
By \cref{eq:nplus-error}, the absolute error tends to zero.  It is therefore smaller than $1/2$ for all sufficiently large $n$, at which point nearest-integer rounding is exact.
\end{proof}

\begin{remark}
The theorem is asymptotic and does not claim an optimal threshold.  Direct computation in \cref{sec:numerics} shows that the dominant-sector inverse is already extraordinarily accurate at modest $n$.
\end{remark}

\subsection{Residual certificates}

An approximate inverse should be checked by substitution rather than by counting formal terms.  Let
\begin{equation}
 F_{K,r}(x)=\log P_{K,r}(x),\qquad
 \rho(\widetilde x)=F_{K,r}(\widetilde x)-\log y.
\end{equation}

\begin{proposition}[A posteriori error certificate]\label{prop:residual-certificate}
Suppose $F_{K,r}'(x)\ge m>0$ on an interval containing the exact root $x_{K,r}(y)$ and an approximation $\widetilde x$.  Then
\begin{equation}\label{eq:residual-bound}
 |\widetilde x-x_{K,r}(y)|\le\frac{|\rho(\widetilde x)|}{m}.
\end{equation}
The corresponding index error satisfies
\begin{equation}\label{eq:residual-index-bound}
 |\widetilde n-n_{K,r}(y)|
 \le \frac{3}{2\pi^2}
 |\widetilde x-x_{K,r}(y)|
 \bigl(|\widetilde x|+|x_{K,r}(y)|\bigr).
\end{equation}
\end{proposition}

\begin{proof}
The mean-value theorem gives
$|F_{K,r}(\widetilde x)-F_{K,r}(x_{K,r})|
=F_{K,r}'(\xi)|\widetilde x-x_{K,r}|$ for an intermediate $\xi$; use the lower bound.  The index estimate is the factorization of the difference of two squares.
\end{proof}

This certificate applies equally to a truncated symbolic transseries and to a numerically refined root.  It is particularly effective because $F_{K,r}'(x)\to1$.

\section{Numerical checks}\label{sec:numerics}

The following computations use exact integer values of $p(n)$, exact rational Dedekind sums inside $A_k(n)$, and high-precision evaluation of the elementary sector formula.  They are intended as verification of normalization and signs, not as substitutes for the proofs.

\subsection{Forward truncations}

Let $P_K(x;n)$ denote \cref{eq:PKr} evaluated with the actual amplitudes $A_k(n)$ and the exact phase $x=\lambda\sqrt{n-1/24}$.  The table shows the signed relative error
$(P_K-p(n))/p(n)$.

\begin{table}[htbp]
\centering
\small
\begin{tabular}{r r r r r r}
\toprule
$n$ & $p(n)$ & $K=1$ & $K=2$ & $K=3$ & $K=5$\\
\midrule
25  & 1,958 & $ 1.0717\,10^{-3}$ & $ 3.0045\,10^{-6}$ & $-6.1015\,10^{-5}$ & $-3.4866\,10^{-6}$\\
50  & 204,226 & $-7.3074\,10^{-5}$ & $ 3.9028\,10^{-6}$ & $ 2.1050\,10^{-7}$ & $-1.1350\,10^{-7}$\\
100 & 190,569,292 & $-1.8220\,10^{-6}$ & $ 8.6838\,10^{-9}$ & $-4.9491\,10^{-9}$ & $ 3.1454\,10^{-10}$\\
250 & 230,793,554,364,681 & $-1.0760\,10^{-9}$ & $ 7.2267\,10^{-13}$ & $ 4.3021\,10^{-14}$ & $ 2.1256\,10^{-15}$\\
500 & 2,300,165,032,574,323,995,027 & $-2.4389\,10^{-13}$ & $1.7537\,10^{-17}$ & $-2.0017\,10^{-19}$ & $-3.5933\,10^{-22}$\\
\bottomrule
\end{tabular}
\caption{Relative error of finite Rademacher sector sums.  Because amplitudes can change sign, accuracy need not improve monotonically at every intermediate $K$, although the convergent sum tends to $p(n)$.}
\label{tab:forward-errors}
\end{table}

\subsection{Inverse truncations}

For $y=p(n)$, let $n_X=1/24+3X^2/(2\pi^2)$ be the carrier inverse, let $n_+$ be the complete dominant-sector inverse, and let $n_K$ be obtained by numerically inverting $P_K(x;n)$ with the true residue data.  The table reports $n_{\mathrm{approx}}-n$.

\begin{table}[htbp]
\centering
\small
\begin{tabular}{r r r r r r}
\toprule
$n$ & carrier $n_X$ & dominant $n_+$ & $K=2$ & $K=3$ & $K=5$\\
\midrule
50   & $-3.5044\,10^{-1}$ & $ 4.5102\,10^{-4}$ & $-2.4089\,10^{-5}$ & $-1.2992\,10^{-6}$ & $ 7.0057\,10^{-7}$\\
100  & $-3.3600\,10^{-1}$ & $ 1.5378\,10^{-5}$ & $-7.3293\,10^{-8}$ & $ 4.1772\,10^{-8}$ & $-2.6548\,10^{-9}$\\
500  & $-3.1768\,10^{-1}$ & $ 4.4042\,10^{-12}$ & $-3.1668\,10^{-16}$ & $ 3.6146\,10^{-18}$ & $ 6.4888\,10^{-21}$\\
1000 & $-3.1356\,10^{-1}$ & $ 4.2960\,10^{-17}$ & $-3.1780\,10^{-23}$ & $ 8.7815\,10^{-26}$ & $-9.9717\,10^{-29}$\\
\bottomrule
\end{tabular}
\caption{Errors in inverse indices at exact ordinates $y=p(n)$.  Restoring the factor $1-1/x$ removes the carrier's order-one bias; the remaining errors follow the exponentially small Rademacher hierarchy.}
\label{tab:inverse-errors}
\end{table}

The parity-sensitive $k=2$ correction explains why the sign of the first exponentially small error alternates.  The very small dominant-sector errors at $n=500$ and $1000$ also illustrate \cref{thm:eventual-rounding}: exact integer recovery requires no knowledge of the residue phase once the error is below one half.

\section{Constructive coefficient algorithms}\label{sec:algorithms}

The formulas above are designed to be executable in a computer algebra system, but they require only formal power-series arithmetic.  This section records deterministic algorithms and their dependency structure.

\subsection{Forward coefficient array}

\begin{algorithm}[Coefficients of a Rademacher sector]\label{alg:forward-sector}
Given $(k,\sigma,M)$:
\begin{enumerate}
\item For $1\le j\le M$, evaluate the finite sum \cref{eq:ell-explicit} to obtain $\ell_{k,\sigma,j}$.
\item Set $c_{k,\sigma,0}=1$.
\item For $1\le m\le M$, use \cref{eq:c-recurrence}:
\[
 c_{k,\sigma,m}=\frac1m\sum_{j=1}^{m}
 j\ell_{k,\sigma,j}c_{k,\sigma,m-j}.
\]
\item Return the sector \cref{eq:sector-unshifted} truncated after $m=M$.
\end{enumerate}
The arithmetic complexity is $\bigO(M^2)$ coefficient operations after the $\ell_j$ are available.  If only the dominant sector is needed, the finite formula \cref{eq:omega-closed} evaluates each $\omega_m$ independently.
\end{algorithm}

\subsection{Lambert carrier and algebraic inverse}

\begin{algorithm}[Phase-independent inverse]\label{alg:dominant-inverse}
Given a large ordinate $y$ and algebraic order $M$:
\begin{enumerate}
\item Compute $X$ from \cref{eq:X-Lambert}; alternatively, compute $R,L$ and use \cref{eq:X-full-log} with $P_j$ from \cref{eq:Pj-stirling}.
\item Generate $b_1,\ldots,b_M$ by \cref{eq:b-recurrence-compact}.
\item Form $x_0=X+\sum_{m=1}^{M}b_mX^{-m}$.
\item Return $n_+=1/24+3x_0^2/(2\pi^2)$.
\item Optionally, evaluate the residual
\[
 \rho=x_0-2\log x_0+\log(1-x_0^{-1})-\log(y/\Ccal)
\]
and apply \cref{prop:residual-certificate}.
\end{enumerate}
For an exact ordinate $y=p(n)$, nearest-integer rounding of the result is eventually exact.
\end{algorithm}

\subsection{Arithmetic exponential correction}

\begin{algorithm}[Finite-level full inverse]\label{alg:full-inverse}
Given $y$, level $K$, compatible residue $r\pmod{L_K}$, and a finite set of desired multi-indices:
\begin{enumerate}
\item Compute $x_0$ by \cref{alg:dominant-inverse} or by numerical solution of \cref{eq:dominant-inverse-equation}.
\item Compute the exact rational Dedekind sums $s(h,k)$ and amplitudes $A_k(r)$ for $2\le k\le K$.
\item Build the gaps and rational amplitudes in \cref{eq:g1minus,eq:gkplus,eq:gkminus}.
\item Traverse the multi-indices in a monomial order and compute $V_\alpha$ by \cref{eq:Valpha-recurrence}.
\item Substitute $q_j=\e^{-\delta_jx_0}$, form $x=x_0+v$, and return $n=1/24+3x^2/(2\pi^2)$.
\item Certify the result by the logarithmic residual of $P_{K,r}(x)-y$ and include the analytic Rademacher tail estimate if comparison with the exact sequence is required.
\end{enumerate}
\end{algorithm}

\begin{remark}[Avoiding the apparent circularity]
For an unknown exact ordinate $y=p(n)$, the residue $r$ is initially unknown.  Compute $n_+(y)$ first and round it.  By \cref{thm:eventual-rounding}, this recovers $n$ for all sufficiently large inputs, after which every desired residue and Rademacher amplitude is known.  Thus the arithmetic transseries refines an already identifiable integer; it is not needed to guess the phase blindly.
\end{remark}

\subsection{Coefficient reuse}

The calculations have a useful modular architecture.
\begin{itemize}
\item The polynomials $P_j(L)$ depend only on the carrier exponent $X^{-2}\e^X$.
\item The rational numbers $b_m$ depend only on the amplitude $1-1/x$.
\item The inverse derivative polynomials $h^{(m)}$ depend only on the dominant phase $\Phi$.
\item The arithmetic data enter only through $A_k(r)$ in the sector amplitudes $g_j$.
\end{itemize}
Consequently, the universal coefficient tables may be precomputed once, while residue-dependent evaluations are inexpensive.

\section{A general template for Rademacher-type sequences}\label{sec:generalization}

The derivation is not limited to ordinary partitions.  Suppose a sequence has an exact or exponentially accurate representation
\begin{equation}\label{eq:general-rademacher-type}
 a(n)=M(x(n))
 \left(1+\sum_{j\in J}c_j(n)g_j(x(n))\e^{-\delta_jx(n)}\right),
\end{equation}
where
\begin{enumerate}
\item $x(n)\to\infty$ is invertible at the algebraic level;
\item $M(x)=C\e^{\alpha x}x^{-\beta}A(x^{-1})$ with $A(0)=1$;
\item $c_j(n)$ is periodic or otherwise belongs to a finite-dimensional phase system at each truncation level;
\item finite subsets of sectors have exponentially smaller tails.
\end{enumerate}
Then the method of this article applies verbatim:
\begin{enumerate}
\item invert $C\e^{\alpha x}x^{-\beta}$ with a Lambert $W$ core;
\item invert the amplitude $A(x^{-1})$ by an ordinary implicit power series;
\item freeze the finite phase data and invert the remaining sectors by \cref{eq:Valpha-recurrence} or \cref{eq:LB-x};
\item convert $x$ back to the original index.
\end{enumerate}

For the general carrier
\begin{equation}
 y=C\e^{\alpha X}X^{-\beta},\qquad \alpha,\beta>0,
\end{equation}
one obtains
\begin{equation}\label{eq:general-carrier-W}
 X=-\frac{\beta}{\alpha}
 W_{-1}\!\left(
 -\frac{\alpha}{\beta}
 \left(\frac{C}{y}\right)^{1/\beta}
 \right).
\end{equation}
The partition case is $(\alpha,\beta)=(1,2)$.  The same Bell-polynomial and triangular coefficient calculus therefore applies to many eta-quotient coefficients and Rademacher-type formulas, provided the relevant phase and tail hypotheses are verified.

\begin{proposition}[General first exponential inverse correction]\label{prop:general-first-correction}
Let $M_0(x)$ be a strictly increasing dominant approximation and suppose
\[
 a=M_0(x)\bigl(1+g(x)\e^{-\delta x}+\bigO(\e^{-2\delta x})\bigr).
\]
If $x_0$ solves $M_0(x_0)=a$, then
\begin{equation}\label{eq:general-first-correction}
 x=x_0-\frac{g(x_0)}{(\log M_0)'(x_0)}\e^{-\delta x_0}
 +\bigO(\e^{-2\delta x_0})
\end{equation}
up to the algebraic factors generated by differentiation.
\end{proposition}

\begin{proof}
Put $x=x_0+v$, take logarithms, and retain the first order in both $v$ and $\e^{-\delta x_0}$.
\end{proof}

This proposition is the one-sector form of \cref{eq:first-v-actual}.  The full multivariate recurrence is needed whenever products, resonances, or several periodic phases interact.

\section{Structural consequences and research directions}\label{sec:research-directions}

\subsection{Convergent perturbative blocks inside an exponentially complete object}

The partition transseries has an uncommon hybrid structure.  The exact Rademacher sum converges.  Within every fixed sector, the unshifted coefficient series \cref{eq:Cksigma} is the Taylor series of an analytic germ and therefore converges near $t=0$.  Likewise, the dominant inverse correction $U(z)$ is analytic near $z=0$.  Divergence is not needed to create the nonperturbative sectors: those sectors are already supplied by modularity and the circle method.

This illustrates a useful distinction:
\begin{itemize}
\item a \emph{Poincare expansion} records a single dominant sector and forgets all exponentially small information;
\item a \emph{transseries completion} records the other exponential sectors, whether or not the perturbative series in each sector diverges;
\item an \emph{exact transseries representation} may itself be convergent, as Rademacher's formula is.
\end{itemize}

\subsection{Arithmetic resurgence without analytic continuation in the index}

The amplitudes $A_k(n)$ are arithmetic rather than Stokes constants.  Nevertheless, they play a formally analogous role: they weight distinct exponential sectors and control exponentially small corrections to inversion.  Their dependence on $n$ is not naturally encoded by a single real analytic function.  The profinite parameter of \cref{thm:profinite} provides a clean replacement: every finite exponential resolution sees only finitely much congruence information.

A promising next step is to formulate an ``arithmetic transseries sheaf'' whose base is $\widehat{\ZZ}$ and whose fibers are finite-level analytic transseries algebras.  The projective compatibility proved here is the elementary case of such a structure.

\subsection{Sharpening inverse error bounds}

The crude estimate \cref{eq:nK-error} uses only $|A_k(n)|\le k$.  Stronger bounds for Rademacher/Kloosterman sums and Lehmer's explicit remainders should yield:
\begin{enumerate}
\item effective constants in the inverse error;
\item a certified minimal threshold for nearest-integer recovery from $n_+(p(n))$;
\item an adaptive rule choosing $K$ from a target number of digits;
\item rigorous interval enclosures for the generalized inverse $\nu(y)$.
\end{enumerate}
The residual method in \cref{prop:residual-certificate} separates the local inversion error from the Rademacher tail, making such a refinement modular.

\subsection{Canonical real interpolation}

The finite-level residue interpolation is canonical along arithmetic subsequences, but it does not produce a unique global real interpolation of $p(n)$.  Possible global choices include:
\begin{itemize}
\item trigonometric interpolation of each periodic $A_k$;
\item logarithmic piecewise-linear interpolation through the points $(n,\log p(n))$;
\item a Newton series or cardinal interpolation satisfying prescribed growth conditions;
\item analytic continuations based on Poincare series or modular integrals.
\end{itemize}
Different choices may have the same algebraic inverse and the same values on the integer range while differing in exponentially small off-lattice terms.  Classifying these differences is a natural nonuniqueness problem.

\subsection{Formal verification targets}

The exact Rademacher input now has an Isabelle/HOL formalization~\cite{EberlPaulson2026}.  The remaining steps are well suited to staged formalization:
\begin{enumerate}
\item differentiation and algebraic normalization to \cref{eq:exact-sector};
\item the coefficient generator \cref{eq:Cksigma} and recurrence \cref{eq:c-recurrence};
\item the Catalan/Lagrange proof of \cref{eq:omega-closed};
\item formal implicit-series construction of $U$ and the recurrence \cref{eq:b-recurrence-compact};
\item finite multivariate triangular inversion \cref{eq:Valpha-recurrence};
\item residual-to-root error transfer under an explicit derivative lower bound.
\end{enumerate}
The Bell-polynomial infrastructure described in the ProveIt combinatorial coefficient calculus is directly relevant to steps 2, 4, and 5.

\section{Conclusion}\label{sec:conclusion}

The ordinary partition numbers admit a particularly explicit transseries description once Rademacher's exact formula is written in the shifted phase $x=\pi\sqrt{2(n-1/24)/3}$.  The main conclusions are:

\begin{enumerate}
\item The full forward transseries is the exact convergent sector sum \cref{eq:exact-sector}.  Every sector has a two-term amplitude in $x^{-1}$.
\item Re-expansion in $n^{-1/2}$ is generated by the universal analytic function \cref{eq:Cksigma}.  Every coefficient is given by the Bell formula \cref{eq:c-bell} or recurrence \cref{eq:c-recurrence}; the dominant coefficients also have the finite formula \cref{eq:omega-closed}.
\item The elementary inverse carrier is exactly a $W_{-1}$ expression.  Its full logarithmic expansion has the Stirling-number coefficient polynomials \cref{eq:Pj-stirling}.
\item Restoring the dominant amplitude produces a convergent rational correction $U(1/X)$ with the all-order recurrence \cref{eq:b-recurrence-compact}.
\item The remaining inverse sectors are arithmetic.  At finite Rademacher level and fixed residue class, every coefficient is generated by the triangular recurrence \cref{eq:Valpha-recurrence}, equivalently by the Lagrange--Buer\-mann/Bell formula \cref{eq:LB-x,eq:vm-Bell}.  Compatible levels are indexed by a profinite phase.
\item For exact ordinates $y=p(n)$, the phase-independent dominant inverse differs from $n$ by an exponentially small quantity.  It therefore recovers the exact integer by rounding for all sufficiently large $n$.
\end{enumerate}

The resulting nested formula \cref{eq:master-inverse} is a complete constructive answer to the asymptotic inversion problem: it separates the universal Lambert/logarithmic geometry, the rational amplitude correction, and the periodic arithmetic sectors, while supplying general coefficient formulas at every layer.

\appendix

\section{Coefficient tables}\label{app:tables}

\subsection{Carrier polynomials}

For convenience, the next two carrier polynomials after those in \cref{eq:P-first} are
\begin{align}
 P_6(L)
 &=-\frac{16}{15}L
 \left(20L^5-274L^4+1125L^3-1700L^2+900L-120\right),\\
 P_7(L)
 &=\frac{32}{105}L
 \left(120L^6-2058L^5+11368L^4-25725L^3
 +24500L^2-8820L+840\right).
\end{align}
They are generated without series guessing by \cref{eq:Pj-stirling}.

\subsection{Dominant inverse coefficients}

The sequence $b_m$ begins
\begin{equation}
 1,\ \frac52,\ \frac{13}{3},\ \frac{53}{12},\ -\frac{14}{5},\
 -\frac{763}{30},\ -\frac{2179}{35},\ -\frac{42289}{630},\
 \frac{132973}{1260},\ \frac{3449711}{5040},\ldots.
\end{equation}
A useful implementation check is the residual identity
\begin{equation}\label{eq:U-residual-order}
 U_M+\log(1-z+zU_M)-3\log(1+zU_M)=\bigO(z^{M+1}),
\end{equation}
where $U_M=\sum_{m=1}^{M}b_mz^m$.

\subsection{First logarithmic sector coefficients}

For arbitrary $k$ and $\sigma$, \cref{eq:ell-explicit} gives
\begin{align}
 \ell_{k,\sigma,1}
 &=-\frac{\sigma\lambda a}{2k}-\frac{\sigma k}{\lambda},\\
 \ell_{k,\sigma,2}
 &=a-\frac{k^2}{2\lambda^2},\\
 \ell_{k,\sigma,3}
 &=-\frac{\sigma\lambda a^2}{8k}
 -\frac{\sigma ka}{2\lambda}
 -\frac{\sigma k^3}{3\lambda^3},\\
 \ell_{k,\sigma,4}
 &=\frac{a^2}{2}
 -\frac{k^2a}{2\lambda^2}
 -\frac{k^4}{4\lambda^4}.
\end{align}
Consequently,
\begin{align}
 c_{k,\sigma,1}&=\ell_{k,\sigma,1},\\
 c_{k,\sigma,2}&=\ell_{k,\sigma,2}+\frac12\ell_{k,\sigma,1}^2,\\
 c_{k,\sigma,3}&=\ell_{k,\sigma,3}
 +\ell_{k,\sigma,1}\ell_{k,\sigma,2}
 +\frac16\ell_{k,\sigma,1}^3,\\
 c_{k,\sigma,4}&=\ell_{k,\sigma,4}
 +\ell_{k,\sigma,1}\ell_{k,\sigma,3}
 +\frac12\ell_{k,\sigma,2}^2
 +\frac12\ell_{k,\sigma,1}^2\ell_{k,\sigma,2}
 +\frac1{24}\ell_{k,\sigma,1}^4.
\end{align}
These are the first instances of \cref{eq:c-bell}.

\section{Derivation of the first inverse perturbation coefficients}\label{app:perturbation}

For completeness, substitute
$x=x_0+\tau v_1+\tau^2v_2+\tau^3v_3+\cdots$ into
\[
 \Phi(x)+\tau\Psi(x)=R,
 \qquad \Phi(x_0)=R.
\]
At orders $\tau,\tau^2,\tau^3$ one obtains
\begin{align*}
 Dv_1+\Psi&=0,\\
 Dv_2+\frac12\Phi''v_1^2+\Psi'v_1&=0,\\
 Dv_3+\Phi''v_1v_2+\frac16\Phi'''v_1^3
 +\Psi'v_2+\frac12\Psi''v_1^2&=0.
\end{align*}
Solving successively yields \cref{eq:v1-diff,eq:v2-diff,eq:v3-diff}.  The same calculation at arbitrary order is encoded compactly by partial Bell polynomials in \cref{eq:vm-Bell}.

\section{Normalization cross-checks}\label{app:normalization}

Several equivalent expressions are useful for detecting missing factors of $k$ or $\sqrt{k}$.
Starting from \cref{eq:rademacher-input}, the $k$th derivative factor is
\begin{align}
 \frac{\dd}{\dd n}\left(N^{-1/2}\sinh(x/k)\right)
 &=\frac1{2N^{3/2}}
 \left(\frac{x}{k}\cosh(x/k)-\sinh(x/k)\right)\\
 &=\frac{x}{4kN^{3/2}}
 \left[\e^{x/k}(1-k/x)+\e^{-x/k}(1+k/x)\right].
\end{align}
After multiplication by $\sqrt{k}/(\pi\sqrt2)$, this may be written either as
\begin{equation}
 \frac1{4\sqrt{3k}\,N}
 \left[\e^{x/k}(1-k/x)+\e^{-x/k}(1+k/x)\right]
\end{equation}
or as
\begin{equation}
 \frac{\Ccal}{\sqrt{k}\,x^2}
 \left[\e^{x/k}(1-k/x)+\e^{-x/k}(1+k/x)\right].
\end{equation}
For $k=1$, the positive term is
\begin{equation}
 \frac{\e^x}{4\sqrt3\,N}\left(1-\frac1x\right),
\end{equation}
which tends to the Hardy--Ramanujan leading expression after replacing $N$ by $n$ and $x$ by $\lambda\sqrt n$.

\section{Reference implementation blueprint}\label{app:implementation}

The following language-neutral pseudocode emphasizes exact coefficient arithmetic.

\begin{lstlisting}[basicstyle=\ttfamily\small,frame=single,breaklines=true]
ForwardSectorCoefficients(k, sigma, M):
    ell[1..M] := formula (5.8)
    c[0] := 1
    for m = 1..M:
        c[m] := (1/m) * Sum_{j=1..m} j*ell[j]*c[m-j]
    return c

DominantInverseCoefficients(M):
    U := 0
    for m = 1..M:
        b[m] := -Coeff(z^m,
             log(1-z+z*U) - 3*log(1+z*U))
        U := U + b[m]*z^m
    return b

ArithmeticInverseCoefficients(K, residue, index_set):
    construct A_k(residue), gaps delta_j, amplitudes g_j
    v := 0
    for alpha in chosen monomial order over index_set:
        rhs := Sum_{m>=2} kappa_m*v^m
             + log(1 + Sum_j q_j*g_j(x0+v)*exp(-delta_j*v))
        V[alpha] := -Coeff(q^alpha, rhs)/D(x0)
        v := v + V[alpha]*q^alpha
    return V
\end{lstlisting}

For numerical work, evaluate logarithms of the dominant terms to avoid overflow, retain exact rational values for Dedekind sums as long as possible, and use the residual certificate of \cref{prop:residual-certificate} after substituting numerical exponentials.

\begin{thebibliography}{99}

\bibitem{Andrews1976}
G.~E. Andrews,
\emph{The Theory of Partitions},
Addison--Wesley, Reading, Massachusetts, 1976; reprinted by Cambridge University Press, 1998.

\bibitem{BanerjeeEtAl2022}
K.~Banerjee, P.~Paule, C.-S. Radu, and C.~Schneider,
``Error bounds for the asymptotic expansion of the partition function,''
arXiv:2209.07887, 2022.
\url{https://arxiv.org/abs/2209.07887}

\bibitem{BrassescoMeyroneinc2020}
S.~Brassesco and A.~Meyroneinc,
``An expansion for the number of partitions of an integer,''
\emph{Ramanujan Journal} \textbf{51} (2020), 563--592.

\bibitem{CorlessEtAl1996}
R.~M. Corless, G.~H. Gonnet, D.~E.~G. Hare, D.~J. Jeffrey, and D.~E. Knuth,
``On the Lambert $W$ function,''
\emph{Advances in Computational Mathematics} \textbf{5} (1996), 329--359.

\bibitem{Comtet1974}
L.~Comtet,
\emph{Advanced Combinatorics: The Art of Finite and Infinite Expansions},
D. Reidel, Dordrecht, 1974.

\bibitem{EberlPaulson2026}
M.~Eberl and L.~C. Paulson,
``Rademacher's Series for the Partition Function,''
Archive of Formal Proofs, 3 July 2026.
\url{https://isa-afp.org/entries/Rademacher_Series.html}

\bibitem{HardyRamanujan1918}
G.~H. Hardy and S.~Ramanujan,
``Asymptotic formulae in combinatory analysis,''
\emph{Proceedings of the London Mathematical Society} (2) \textbf{17} (1918), 75--115.

\bibitem{Kong2023}
Z.~Y. Kong,
``Rademacher's formula for the partition function,''
arXiv:2302.03835, 2023.
\url{https://arxiv.org/abs/2302.03835}

\bibitem{Lehmer1937}
D.~H. Lehmer,
``On the Hardy--Ramanujan series for the partition function,''
\emph{Journal of the London Mathematical Society} \textbf{12} (1937), 171--176.

\bibitem{Lehmer1938}
D.~H. Lehmer,
``On the series for the partition function,''
\emph{Transactions of the American Mathematical Society} \textbf{43} (1938), 271--295.

\bibitem{Lehmer1939}
D.~H. Lehmer,
``On the remainders and convergence of the series for the partition function,''
\emph{Transactions of the American Mathematical Society} \textbf{46} (1939), 362--373.

\bibitem{OEISA000041}
OEIS Foundation Inc.,
``A000041: Number of partitions of $n$ (the partition numbers),''
\url{https://oeis.org/A000041}, accessed 3 September 2026.

\bibitem{OSullivan2023}
C.~O'Sullivan,
``Detailed asymptotic expansions for partitions into powers,''
\emph{International Journal of Number Theory} \textbf{19} (2023), no.~9, 2163--2196;
arXiv:2205.13468.
\url{https://arxiv.org/abs/2205.13468}

\bibitem{ProveItCCC}
V.~Reshetnikov,
``Combinatorial Coefficient Calculus,''
ProveIt repository, 2026.
\url{https://github.com/VladimirReshetnikov/ProveIt/tree/main/Analysis/FabiusFunction/docs/semi-formalized-research-frontiers/drafts/combinatorial-coefficient-calculus/Combinatorial_Coefficient_Calculus}

\bibitem{ProveItTransseries}
V.~Reshetnikov,
``Series and Transseries research drafts,''
ProveIt repository, 2026.
\url{https://github.com/VladimirReshetnikov/ProveIt/tree/main/Analysis/FabiusFunction/docs/semi-formalized-research-frontiers/drafts/series-and-transseries}

\bibitem{Rademacher1937}
H.~Rademacher,
``A convergent series for the partition function $p(n)$,''
\emph{Proceedings of the National Academy of Sciences USA} \textbf{23} (1937), 78--84.

\bibitem{Rademacher1938}
H.~Rademacher,
``On the partition function $p(n)$,''
\emph{Proceedings of the London Mathematical Society} (2) \textbf{43} (1938), 241--254.

\bibitem{Rademacher1943}
H.~Rademacher,
``On the expansion of the partition function in a series,''
\emph{Annals of Mathematics} \textbf{44} (1943), 416--422.

\bibitem{Szekeres1953}
G.~Szekeres,
``Some asymptotic formulae in the theory of partitions (II),''
\emph{Quarterly Journal of Mathematics} \textbf{4} (1953), 96--111.

\bibitem{Wright1934}
E.~M. Wright,
``Asymptotic partition formulae III: partitions into $k$th powers,''
\emph{Acta Mathematica} \textbf{63} (1934), 143--191.

\end{thebibliography}

\end{document}
